\chapter{Le groupe fondamental: g\'en\'eralit\'es}
\label{V}

\setcounter{section}{-1}
\section{Introduction}

Le pr\'esent S\'eminaire est la suite du S\'eminaire 1960. Nous
r\'ef\'erons \`a ce dernier par des sigles tels que I~\ref{I.9.7}
qui signifie: S\'eminaire de G\'eom\'etrie Alg\'ebrique,
expos\'e~I, \No9.7. Les num\'eros des expos\'es de 1961 suivront
ceux de 1960. Nous r\'ef\'erons aux \'El\'ements de
G\'eom\'etrie Alg\'ebrique de Dieudonn\'e-Grothendieck par des
sigles tels que (EGA~I 8.7.3).

Le pr\'esent expos\'e r\'esume (avec de l\'egers
compl\'ements) les derniers expos\'es de 1960, qui n'avaient pas
\'et\'e r\'edig\'es.

Comme en 1961, nous nous limiterons en r\`egle g\'en\'erale
\`a des pr\'esch\'emas localement noeth\'eriens, bien que
souvent cette restriction soit inessentielle. Nous admettrons dans
l'expos\'e \ref{VI} la th\'eorie de la descente fid\`element plate,
r\'esum\'ee dans S\'eminaire Bourbaki \No190. S'il y a lieu,
nous en donnerons un expos\'e plus d\'etaill\'e dans un
expos\'e ult\'erieur\footnote{\Cf Exp.\, \ref{VI} et Exp.\, \ref{VIII}}, une fois
que le lecteur aura eu l'occasion de se convaincre de l'utilit\'e de
cette technique, pour la th\'eorie du groupe fondamental.

\section{Pr\'esch\'ema \`a groupe fini d'op\'erateurs,
pr\'esch\'ema quotient}
\label{V.1}

Soient $X$ un pr\'esch\'ema, $G$ un groupe fini op\'erant sur~$X$ par automorphismes, \`a droite pour fixer les id\'ees. Si $X$
est affine d'anneau $A$, $G$ op\`ere donc par automorphismes \`a
gauche sur~$A$.

Pour tout pr\'esch\'ema $Z$, $G$ op\`ere \`a gauche sur
l'ensemble $\Hom(X,Z)$,
on peut donc consid\'erer l'ensemble
$$
\Hom(X,Z)^G
$$
des morphismes invariants par $G$. Il d\'epend fonctoriellement de
$Z$,
on peut se demander si ce foncteur est \og repr\'esentable\fg,
\ie isomorphe \`a un foncteur $Z\mto\Hom(Y,Z)$.
Cela signifie qu'on peut trouver un pr\'esch\'ema $Y$, et un
morphisme invariant par $G$
$$
p\colon X\to Y
$$
tel que pour tout $Z$, l'application correspondante $g\mto gp$
$$
\Hom(Y,Z)\to \Hom(X,Z)^G
$$
soit bijective. On dit alors que $(Y,p)$ est un
\emph{préschéma quotient}
\index{quotient (pr\'esch\'ema)|hyperpage}de~$X$
par~$G$
(il est d\'etermin\'e \`a isomorphisme unique pr\`es).

\begin{proposition}
\label{V.1.1}
Soient $A$ un anneau sur lequel le groupe fini $G$ op\`ere \`a
gauche, $B=A^G$ le sous-anneau des invariants de~$A$, $X=\Spec(A)$ et
$Y=\Spec(B)$, $p\colon X\to Y$ le morphisme canonique (\'evidemment
invariant par $G$). Alors
\begin{enumerate}
\item[(i)] $A$ est entier sur~$B$ \ie $p$ est un morphisme
\emph{entier}.
\item[(ii)] Le morphisme $p$ est surjectif, ses fibres sont les
trajectoires de~$G$, la topologie de~$Y$ est quotient de celle de~$X$.
\item[(iii)] Soit $x\in X$, $y=p(x)$, $G_x$ le stabilisateur de~$x$, alors
$\kres(x)$ est une extension alg\'ebrique quasi-galoisienne de
$\kres(y)$ et l'application canonique de~$G_x$ dans le groupe
$\Gal(\kres(x)/\kres(y))$ des $\kres(y)$-automorphismes de
$\kres(x)$
est
surjective.
\item[(iv)] $(Y,p)$ est un pr\'esch\'ema quotient de~$X$ par $G$.
\end{enumerate}
\end{proposition}

Les \'enonc\'es
(i), (ii), (iii)
sont bien connus en alg\`ebre
commutative\footnote{\Cf N. Bourbaki, Alg. Comm. Chap.\, 5, \S 1 et \S
2, th.\, 2.} et sont mis seulement pour m\'emoire, sauf l'assertion
sur la topologie, qui provient du fait g\'en\'eral suivant,
cons\'equence facile du th\'eor\`eme de Cohen-Seidenberg: un
morphisme entier est ferm\'e (\ie transforme ferm\'es en
ferm\'es). Notons tout de suite:

\begin{corollaire}
\label{V.1.2}
Sous les conditions pr\'ec\'edentes, l'homomorphisme naturel
$$\cal{O}_Y{\to} p_*(\cal{O}_X)^G$$ est un isomorphisme.
\end{corollaire}

Cela r\'esulte aussit\^ot de la formule
$$
(S^{-1}A)^G = S^{-1}(A^G)
$$
valable pour toute partie multiplicativement stable $S$ de~$B=A^G$
(formule qui se module, et s'\'enonce plus g\'en\'eralement pour
un changement de base $A\to A'$ qui est \emph{plat}),
appliqu\'ee
au cas o\`u $S$ est engendr\'e par un \'el\'ement $f$ de~$B$.

L'assertion (ii) et cor.\, \ref{V.1.2} impliquent facilement (iv); plus
g\'en\'eralement, on aura ceci:

\begin{proposition}
\label{V.1.3}
Soient
$X$ un pr\'esch\'ema \`a groupe d'automorphismes finis
$G$, $p\colon{}X\to{}Y$ un morphisme affine invariant tel que
$\cal{O}_Y\isomto p_*(\cal{O}_X)^G$. Alors les conclusions
\textup{(i), (ii), (iii), (iv)}
de~\ref{V.1.1} sont encore valables.
\end{proposition}

En effet, pour
(i), (ii), (iii),
on peut supposer $Y$ donc $X$ affine, et
si $B,A$ sont leurs anneaux, l'hypoth\`ese implique $B=A^G$, il
suffit d'appliquer~\ref{V.1.1}. Pour (iv), on utilise (ii) et
$\cal{O}_Y=p_*(\cal{O}_X)^G$.

\begin{corollaire}
\label{V.1.4}
Sous les conditions de~\ref{V.1.3}, pour tout ouvert $U$ de~$Y$, $U$
est un quotient de~$X|U=p^{-1}(U)$ par $G$.
\end{corollaire}

En effet, $p^{-1}(U)\to{}U$ induit par $p$ satisfait aux m\^emes
hypoth\`eses que~$p$.

Si maintenant $X$ est un $Z$-pr\'esch\'ema et les op\'erations
de~$G$ sont des $Z$-automorphismes, alors par (iv) $Y$ est un
$Z$-pr\'esch\'ema. Ceci dit:

\begin{corollaire}
\label{V.1.5}
Pour que $X$ soit affine \resp s\'epar\'e sur~$Z$, il faut et il
suffit que~$Y$ le soit. Si $X$ est de type fini sur~$Z$, il est
\emph{fini} sur~$Y$; si de plus $Z$ est localement noeth\'erien, $Y$
est de type fini sur~$Z$.
\end{corollaire}

Comme $X$ est affine et a fortiori s\'epar\'e sur~$Y$, si $Y$ est
affine \resp s\'epar\'e sur~$Z$, $X$~l'est
aussi. R\'eciproquement, supposons $X$ affine sur~$Z$, prouvons que
$Y$ l'est: on peut gr\^ace \`a~\ref{V.1.4} supposer $Z$ affine, et
on est ramen\'e \`a prouver que si $X$ est affine, $Y$~l'est, ce
qui r\'esulte de la d\'etermination explicite de~$Y$ comme
$\Spec(A)^G$ faite dans~\ref{V.1.1}. De m\^eme, comme
$p\colon{}X\to{}Y$ est entier donc universellement ferm\'e, et
surjectif, il s'ensuit que si $X$ est s\'epar\'e sur~$Z$, $Y$
l'est aussi (lemme \`a d\'egager!); en effet, dans le diagramme
$$
\xymatrix@C=1.5cm{{X \times_ZX}\ar[r]^-{p\times_Zp}&{Y\times_ZY}\\{X}\ar[u]^{\Delta_{X/Z}}\ar[r]^{p}&{Y}\ar[u]_{\Delta_{Y/Z}}}
$$
le morphisme $X\times_ZX\to Y\times_ZY$ est ferm\'e, donc transforme
la diagonale (ferm\'ee) de~$X\times_ZX$ en une partie ferm\'ee de
$Y\times_ZY$, qui n'est d'ailleurs autre que la diagonale de ce
dernier puisque $p$ est surjectif.
--- Si $X$ est de type fini sur~$Z$, il l'est a fortiori sur~$Y$ donc
il est fini sur~$Y$ (puisqu'il est d\'ej\`a entier sur~$Y$). Supposons de plus $Z$ localement noeth\'erien, prouvons que
$Y$ est de type fini sur~$Z$. On peut gr\^ace \`a~\ref{V.1.4}
supposer $Z$ affine. Comme l'espace topologique $X$ est quasi-compact
et que $p\colon X\to Y$ est surjectif, $Y$ est
\'egalement
quasi-compact donc r\'eunion finie d'ouverts
affines, et par~\ref{V.1.4} on est ramen\'e au cas o\`u $Y$ est
affine, donc $X$ affine. Mais alors l'anneau $A$ de~$X$ est une
alg\`ebre de type fini sur l'anneau $C$ de~$Z$ qui est
noeth\'erien, et il est connu que $B=A^G$ est alors \'egalement
une alg\`ebre de type fini sur~$C$ (car $A$ sera enti\`ere, donc
finie sur une sous-alg\`ebre $B'$ de~$B$ de type fini sur~$C$, donc
comme $B'$ est noeth\'erien, $B$ est \'egalement
finie
sur~$B'$, donc de type fini sur~$C$).

\begin{corollaire}
\label{V.1.6}
Pour que $X$ soit affine \resp un sch\'ema, il \fets que $Y$ le
soit.
\end{corollaire}

\begin{definition}
\label{V.1.7}
Soit
$X$ un pr\'esch\'ema o\`u un groupe fini $G$ op\`ere
\`a droite. On dit que $G$ \emph{opère de fa\c con admissible}
s'il existe un morphisme $p\colon X\to Y$ ayant les propri\'et\'es
de~\ref{V.1.3} (ce qui implique que $X/G$ existe et est isomorphe
\`a~$Y$).
\end{definition}

\begin{proposition}
\label{V.1.8}
Soit $X$ un pr\'esch\'ema o\`u le groupe fini $G$ op\`ere
\`a droite. Pour que $G$ op\`ere de fa\c con admissible, il faut
et il suffit que $X$ soit r\'eunion d'ouverts affines invariants par
$G$, ou encore que toute trajectoire de~$G$ dans $X$ soit contenue
dans un ouvert affine.
\end{proposition}

Cette derni\`ere condition est \'evidemment impliqu\'ee par la
premi\`ere, et \`a son tour elle l'implique; car soit $T$ une
trajectoire de~$G$, $U$ un ouvert affine la contenant, l'intersection
des transform\'es de~$U$ par les
$g$ dans $G$
est alors un ouvert $U'$ stable par~$G$, contenant $T$ et contenu dans
l'ouvert affine $U$. Comme dans $U$, toute partie finie a un
syst\`eme fondamental de voisinages ouverts affines, il existe un
voisinage ouvert affine $V$ de~$T$ contenu dans $U'$. Ces
transform\'es par les
$g$ dans $G$
sont donc affines et contenus dans $U'$ qui est \emph{séparé},
donc leur intersection $U''$ est un ouvert affine qui est invariant
par $G$ et contient $T$.
--- Ceci pos\'e, la condition envisag\'ee dans~\ref{V.1.8} est
\emph{nécessaire}, car on prendra les images inverses $X_i$
d'ouverts affines $Y_i$ recouvrant~$Y$. Elle est suffisante, car on
peut alors par~\ref{V.1.1} construire les quotients $Y_i=X_i/G$; dans
chaque $Y_i$ l'image de~$X_i\cap X_j$ est un ouvert $Y_{ij}$
s'identifiant \`a $X_{ij}/G$ par~\ref{V.1.4}, en particulier on en
d\'eduit des isomorphismes $Y_{ij}\isomto Y_{ji}$ permettant de
recoller les $Y_i$ pour construire $Y$. Serre pr\'ef\`ere
construire directement l'espace topologique quotient $Y$ de~$X$ par
$G$, mettre dessus le faisceau $p_*(\cal{O}_X)^G$ et v\'erifier que
$Y$ devient un pr\'esch\'ema et qu'on est alors sous les
conditions de~\ref{V.1.3}.


\setcounter{subsection}{6}

\begin{corollaire}
\label{cor:V.1.7}
Si $G$ op\'erant sur~$X$ est admissible, il en est de m\^eme pour
tout sous-groupe $H$ de~$G$ (donc $X/H$ existe).
\end{corollaire}

Cela peut aussi se v\'erifier directement sur la
situation~\ref{V.1.3}, en notant qu'on peut toujours supposer $X$
affine sur un $Z$ et les $s\in G$ op\`erent par $Z$-automorphismes
(on prend par exemple $Z=Y$); on a en effet:

\begin{corollaire}
\label{cor:V.1.8}
Supposons $X$ affine sur~$Z$, et les op\'erations de~$G$ des
$Z$-auto\-mor\-phismes. Alors $G$ op\`ere sur~$X$ de fa\c con
admissible. Si $X$ est d\'efini par un faisceau quasi-coh\'erent
$\cal{A}$ d'alg\`ebres, $Y$ est d\'efini par le faisceau
$\cal{A}^G$ des invariants de~$\cal{A}$ par~$G$.
\end{corollaire}

\begin{proposition}
\label{V.1.9}
Supposons que $G$ op\`ere de fa\c con admissible sur~$X$, et que
$X/G=Y$ soit un pr\'esch\'ema sur~$Z$. Consid\'erons un
morphisme de changement de base $Z'\to Z$, posons $X'=X\times_Z Z',
Y'=Y\times_Z Z'$, de sorte que $G$ op\`ere encore par transport de
structure sur~$X'$, le morphisme $p'\colon X'\to Y'$ \'etant
invariant. Si $Z'$ est \emph{plat} sur~$Z$, alors $p'$ satisfait
encore les hypoth\`eses
de~\ref{V.1.3},
\ie $\cal{O}_Y'\to
p_*'(\cal{O}_{X'})^G$ est un isomorphisme ($p'$ \'etant de toutes
fa\c cons affine). Donc $G$ op\`ere de fa\c con admissible sur~$X'$, et $(X/G)\times_Z Z'\thickapprox (X\times_Z Z')/G$.
\end{proposition}

On peut \'evidemment supposer $Z=Y$, on est ramen\'e au cas o\`u
de plus $Y$ et $Y'$ sont affines. Il faut montrer que si $B$ est le
sous-anneau des invariants de~$G$ op\'erant dans $A$, et si $B'$ est
une alg\`ebre sur~$B$ plate sur~$B$, alors $B'$ est la
sous-alg\`ebre des invariants de~$A'=A\otimes_B B'$. C'est
imm\'ediat, car la suite exacte
\enlargethispage{.5cm}$$
0 \to B \lto{i} A \lto{j} A^{(G)}
$$
(o\`u le dernier terme signifie une puissance de~$A$, et o\`u
$j(x)$ est le syst\`eme des $s\cdot x-x$, $s\in G$) reste exacte par
tensorisation par
$B'$.

On fera attention que l'hypoth\`ese de platitude \'etait
essentielle pour la validit\'e du r\'esultat; en particulier, si
$Y'$ est un sous-pr\'esch\'ema ferm\'e de
$Y$
(par exemple m\^eme un point ferm\'e de
$Y$),
$X'$ son image inverse dans $X$,
alors $Y'$ \emph{ne s'identifie pas} en g\'en\'eral \`a
$X'/G$. Nous verrons qu'il en est n\'eanmoins ainsi si $X$ est
\'etale sur~$Y$.

Pour finir, donnons un formalisme aussi commode que trivial. Soit $Y$
un pr\'esch\'ema. Comme dans la cat\'egorie des
pr\'esch\'emas, les sommes directes existent, on peut pour tout
ensemble $E$ consid\'erer le pr\'esch\'ema somme d'une famille
$(Y_i)_{i\in E}$ de pr\'esch\'emas tous identiques \`a $Y$, ce
pr\'esch\'ema sera not\'e $Y\times E$.
Il est caract\'eris\'e par la formule
\begin{equation*}
\label{eq:V.1.*}
\tag{$*$} \Hom(Y\times E,Z) = \Hom(E,\Hom(Y,Z))
\end{equation*}
o\`u le deuxi\`eme $\Hom$ d\'esigne \'evidemment l'ensemble
des applications de l'ensemble $E$ dans l'ensemble $\Hom(Y,Z)$. On a
un morphisme canonique
$$
Y\times E\to Y
$$
faisant de~$Y\times E$ un pr\'esch\'ema sur~$Y$. Comme les
produits fibr\'es commutent aux sommes directes (dans la
cat\'egorie des pr\'esch\'emas) on aura, si $Y$ est un
pr\'esch\'ema sur un autre $Z$, pour un changement de base $Z'\to
Z$:
$$
(Y\times E)\times_Z Z'=(Y\times_Z Z')\times E
$$
(formule surtout utile si $Z=Y$). D'autre part, on conclut
trivialement de la d\'efinition
$$
(Y\times E)\times F=Y\times(E\times F)=(Y\times E)\times_Y(Y\times F)
$$
(la derni\`ere formule cependant r\'esultant de la
commutativit\'e signal\'ee plus haut).

Pour $Y$ fix\'e, on peut regarder $Y\times E$ comme un foncteur en
$E$, \`a valeurs dans les pr\'esch\'emas sur~$Y$, foncteur qui
commute aux produits finis d'apr\`es la formule pr\'ec\'edente,
(ce qui permet par exemple \`a tout groupe ordinaire $G$ de faire
correspondre un sch\'ema en groupes $Y\times G$ sur~$Y$, qui sera
fini sur~$Y$ si $Y$ l'est, etc.). Plus g\'en\'eralement, ce
foncteur est \og exact \`a gauche\fg, mais nous n'aurons pas \`a nous
en servir ici. Ce foncteur commute aussi trivialement aux sommes
directes, et il est aussi \og exact \`a droite\fg, comme on voit
aussit\^ot sur la formule de d\'efinition \eqref{eq:V.1.*}. En
particulier, si le groupe fini $G$ op\`ere \`a droite dans
l'ensemble $E$, alors il op\`ere \`a droite dans $Y\times E$, et
on a
$$
(Y\times E)/G = Y\times (E/G)
$$
o\`u en fait le quotient du premier membre satisfait aux conditions
de~\ref{V.1.3} (c'est imm\'ediat).

\section{Groupes de d\'ecomposition et d'inertie. Cas \'etale}
\label{V.2}
Soit $G$
un
groupe fini op\'erant \`a droite sur le pr\'esch\'ema $X$. Si
$x\in X$, on appelle
\emph{groupe de décomposition de~$x$}
\index{decomposition (groupe de)@d\'ecomposition (groupe de)|hyperpage}\index{groupe de d\'ecomposition|hyperpage}le stabilisateur $G_d(x)$ de~$x$. Ce groupe op\`ere canoniquement
(\`a gauche) sur le corps r\'esiduel
$\kres(x)$, et l'ensemble
des \'el\'ements de~$G_d(x)$ qui op\`erent trivialement est
appel\'e \emph{groupe d'inertie}
\index{groupe d'inertie|hyperpage}\index{inertie (groupe d')|hyperpage}de~$x$, not\'e~$G_i(x)$.

Supposons que $G$ op\`ere sur~$X$ de fa\c con admissible et que
$Y$ soit un pr\'esch\'ema sur un pr\'esch\'ema
$Z$. Fixons-nous un $z\in Z$, et une extension alg\'ebriquement
close $\Omega$
de~$\kres(z)$ ayant un degr\'e de transcendance
sup\'erieur \`a celui des $\kres(x)/\kres(z)$, o\`u $x$ est un
point de~$X$ au-dessus de~$z$. On peut regarder $\Spec(\Omega)$ comme
un $Z$-sch\'ema, et les points de~$X$ \`a valeurs dans $\Omega$
correspondent aux homomorphismes
de~$\kres(z)$-alg\`ebres $\kres(x)\to\Omega$,
o\`u $x$ est un point de~$X$ au-dessus de~$z$;
comme $\Omega$ a \'et\'e prise assez grande, tout point $x$ de~$X$
au-dessus de~$z$ est la localit\'e d'un point de~$X$ \`a valeurs
dans $\Omega$. Si $X(\Omega)$ et $Y(\Omega)$ d\'esignent
respectivement l'ensemble des points de~$X$ et $Y$ \`a valeurs dans
$\Omega$, on a une application naturelle
$$
X(\Omega)\to Y(\Omega),
$$
d'autre part, $G$ op\`ere sur~$X(\Omega)$ et l'application
pr\'ec\'edente est invariante par $G$. Ceci pos\'e, les
conclusions (ii) et (iii) de~\ref{V.1.3} s'interpr\`etent aussi
ainsi: \emph{l'application précédente est surjective et
identifie $Y(\Omega)$ au quotient $X(\Omega)/G$. De plus, si $x$ est
la localité de~$a\in X(\Omega)$, alors le stabilisateur de~$a$
dans $G$ n'est autre que le groupe d'inertie $G_i(x)$}. Tout ceci est
d'ailleurs vrai sans supposer $\Omega$ \og assez grand\fg, cette
derni\`ere hypoth\`ese sert uniquement \`a assurer qu'on peut
caract\'eriser le groupe d'inertie de tout \'el\'ement de~$X$
au-dessus de~$z$ comme un stabilisateur \og g\'eom\'etrique\fg. On en
conclut par exemple aussit\^ot:
\begin{proposition}
\label{V.2.1}
Faisons une extension de la base $Z'\to Z$, d'o\`u
$X'=X\times_ZZ'$. Soit $x'$ un point de~$X'$, $x$ son image dans $X$,
alors on a $G_i(x)=G_i(x')$.
\end{proposition}

Il suffit, dans les consid\'erations ci-dessus, de prendre pour
$\Omega$ une extension assez grande
de~$\kres(z')$
(o\`u $z,z'$ sont les images de~$x,x'$ dans $Z,Z'$).

\begin{proposition}
\label{V.2.2}
Sous les conditions de~\ref{V.1.3}, supposons $Y$ localement
noeth\'erien, $X$ fini sur~$Y$. Soit $H$ un sous-groupe de~$G$,
consid\'erons $X'=X/H$ (\cf \ref{cor:V.1.7}), soit
$x\in X$, $x'$ son image dans $X'$ et $y$ son image dans $Y$.
\begin{enumerate}
\item[(i)] Si $H\supset G_d(x)$, alors l'homomorphisme
$\cal{O}_y\to\cal{O}_{x'}$ induit un isomorphisme sur les
compl\'et\'es.
\item[(ii)] Si $H\supset G_i(x)$, alors l'homomorphisme $\cal{O}_y\to
\cal{O}_{x'}$ est \'etale \ie $X'$ est \'etale sur~$Y$ en $x'$.
\end{enumerate}
\end{proposition}
Soit $Y_1=\Spec(\widehat{\cal{O}_y})$, faisons le changement de base
$Y_1\to Y$, on trouve un $X_1=X\times_Y Y_1$ fini sur~$Y_1$, sur
lequel $G$ op\`ere, le quotient \'etant $Y_1$
par~\ref{V.1.9}. Soit~$y_1$ l'unique point de~$Y_1$ au-dessus de~$y$,
comme
$\kres(y)=\kres(y_1)$,
il s'ensuit que la fibre de~$X$ en~$y$
est isomorphe \`a celle de~$X_1$ en $y_1$, d'o\`u un unique point
$x_1$ de~$X_1$ au-dessus de~$x$. D'ailleurs par~\ref{V.1.9} on aura
$X_1/H=X'_1=X'\times_Y Y_1$, soit $x'_1$ l'image de~$x_1$ dans $X'_1$,
il est au-dessus de~$x'$, et on v\'erifie facilement ($X'$ \'etant
de type fini sur~$Y$) que l'homomorphisme
$\cal{O}_{x'}\to\cal{O}_{x'_1}$ induit un isomorphisme sur les
compl\'et\'es. Donc on est ramen\'e au cas o\`u $Y$ est le
spectre d'un anneau local complet, soit $B$, donc~$X$ le spectre d'un
anneau fini $A$ sur~$B$, compos\'e d'un nombre fini d'anneaux locaux~$A_x$ correspondant
aux points $x_i$ de~$X$ sur~$Y$. Si $A_0$ correspond \`a $x=x_0$,
alors $A$ s'identifie \`a l'anneau $\Hom_{G_d}(G,A_0)$ des fonctions
$f\colon G\to A_0$ telles que $f(st)=sf(t)$ pour $s\in G_d$, les
op\'erations de~$G$ sur ces fonctions \'etant d\'efinies par
$(uf)(t)=f(tu)$. On voit donc que si $H$ est un sous-groupe quelconque
de~$G$, alors $A^H$ est l'anneau des fonctions $f\colon G\to A_0$
telles que
$$
f(stu)=sf(t)\qquad\text{pour}\ s\in G_d,\ u\in H
$$
donc c'est un anneau semi-local dont les composants locaux
correspondent aux doubles classes $G_daH$ dans $G$, \`a la double
classe d\'efinie par $a\in G$ correspondant (gr\^ace \`a
l'application $f\mto f(a)$) le sous-anneau $A_0^{H(a)}$ de~$A_0$,
o\`u $H(a)=G_d\cap aHa^{-1}$. D'ailleurs, le composant local de
$A^H$ correspondant \`a l'image $x'$ de~$x$ est aussi celui
correspondant \`a la double classe $G_dH$ de l'\'el\'ement
neutre, son composant local est donc $A_0^{G_d\cap H}$. Si donc
$G_d\subset H$,
on trouve $A_0^{G_d}=A^G=B$, ce qui prouve (i). Pour
prouver (ii), on peut, en passant \`a une extension finie plate
convenable de~$A$, et utilisant~\ref{V.2.1}, se ramener au cas o\`u
l'extension r\'esiduelle
$\kres(x)/\kres(y)$
est triviale. Mais
alors $G_i(x)=G_d(x)$, et on est ramen\'e au cas pr\'ec\'edent.

\begin{corollaire}
\label{V.2.3}
Sous
les conditions de~\ref{V.2.2}, supposons $G_i(x)=(e)$, alors X
est \'etale sur~$Y$ en $x$. Donc si $G_i(x)=(e)$ pour tout $x \in
X$, alors $X \to Y$ est un morphisme \'etale.
\end{corollaire}

Il y a une r\'eciproque partielle:

\begin{corollaire}
\label{V.2.4} Supposons $X$ connexe et le groupe $G$ fid\`ele
sur~$X$. Pour que $p \colon X \to Y=X/G$ soit \'etale, il faut et il
suffit que les groupes d'inertie des points de~$X$ soient r\'eduits
\`a l'\'el\'ement neutre. S'il en est ainsi, $G$ s'identifie au
groupe de tous les $Y$-automorphismes du $Y$-sch\'ema~$X$.
\end{corollaire}

Compte tenu de~\ref{V.2.3}, on peut supposer $X$ \'etale sur~$Y$. Mais si un $s \in G$ est dans un $G_i(x)$, il r\'esulte alors
de I~\ref{I.5.4} que $s$ op\`ere trivialement sur~$G$, donc est
l'\'el\'ement unit\'e puisque $G$ est fid\`ele, ce qui prouve
la premi\`ere assertion. Soit $u$ un $Y$-automorphisme de~$X$, soit
$x \in X$. D'apr\`es la proposition~\ref{V.1.3}, il existe un $s \in
G$ tel que $s(x)=u(x)$, et induisant le m\^eme homomorphisme
r\'esiduel
$\kres(x)/\kres(x')$
que $u$. Par \loccit on a donc
$s=u$, ce qui ach\`eve la d\'emonstration.

\begin{remarque}
\label{V.2.5}
L'hypoth\`ese que $G$ op\`ere fid\`element n'est \'evidemment
par surabondante dans le corollaire~\ref{V.2.4}. Il en est de m\^eme
de l'hypoth\`ese que $X$ est connexe, comme on voit par exemple en
prenant $X=Y \times E$, $E$ \'etant un ensemble fini, et $G$ le
groupe des permutations de~$E$: $G$ op\`ere avec force inertie,
n\'eanmoins $(Y \times E)/G=Y \times (E/G)=Y$, et $X$ est \'etale
sur~$Y$. Prenant pour $G$ un groupe strictement plus petit que le
groupe sym\'etrique de~$E$, mais op\'erant transitivement sur~$E$,
on voit qu'il y aura aussi des $Y$-automorphismes de~$X$ ne provenant
pas de~$G$.
\end{remarque}

L'exemple type d'un groupe $G$ op\'erant sans inertie est celui de
$Y \times G$, sur lequel on fait op\'erer $G$ gr\^ace \`a ses
op\'erations sur le facteur $G$ par translations \`a droite: un
$Y$-pr\'esch\'ema $X$ \`a groupe d'op\'erateurs \`a droite
$G$ est dit \emph{trivial} s'il est isomorphe \`a $Y \times G$.

Pour faire le lien entre le pr\'esch\'emas \`a groupes finis
d'op\'erateurs et la notion de fibr\'e principal dans une
cat\'egorie (lien dont nous n'aurons pas besoin d'ailleurs pour la
suite du s\'eminaire, mais important dans d'autres contextes) les
consid\'erations suivantes sont utiles. Nous fixons un
pr\'esch\'ema de base $Y$, et nous pla\c cons dans le
cat\'egorie des $Y$-pr\'esch\'emas. Si $G$ est un groupe
fini, nous poserons pour abr\'eger $G_Y=Y \times G$, c'est donc un
sch\'ema en groupes finis sur~$Y$ (\cf \numero 1), et si $X$ est
un $Y$-pr\'esch\'ema, on a
$$
X \times_Y G_Y = X \times G
$$
(m\^eme r\'ef\'erence). La donn\'ee d'un $Y$-morphisme $X
\times_Y G_Y \to X$ \'equivaut donc \`a la donn\'ee d'un
$Y$-morphisme $X \times G \to X$, \ie \`a la donn\'ee pour tout
$g \in G$ d'un $Y$-morphisme $T_g \colon X \to X$. On constate
aussit\^ot que pour que la donn\'ee des $T_g$ d\'efinisse sur~$X$ une structure de pr\'esch\'ema \`a groupe d'op\'erateurs
\`a droite $G$, (\ie $T_{gg'}=T_{g'}T_g$, $T_e=\id_x$) il faut et il
suffit que le $Y$-morphisme correspondant $X \times_Y G_Y \to X$
d\'efinisse sur~$X$ une structure de~$Y$-pr\'esch\'ema \`a
$Y$-sch\'ema en groupes d'op\'erateurs, (au sens g\'en\'eral
des objets \`a $\cal{C}$-groupe d'op\'erateurs dans une
cat\'egorie $\cal{C}$). Supposons qu'il en soit ainsi. Rappelons que
$X$ est dit \emph{formellement principal homogène}
\index{formellement principal homog\`ene (pr\'esch\'ema)|hyperpage}sous~$G_Y$\footnote{on dit plut\^ot maintenant: $X$ est un
pseudo-torseur sous $G_Y$.} si le morphisme canonique
$$
X \times_Y G_Y \to X \times_Y X
$$
dont les composantes sont respectivement $\pr_1$ et le morphisme de
multiplication $\pi \colon X \times_Y G_Y \to X$, est un
isomorphisme. En l'occurrence, identifiant le premier membre \`a $X
\times G$, le morphisme consid\'er\'e est celui qui, \`a tout $g
\in G$, associe le morphisme
$$
(\id_X,T_g)=(\id_X \times_Y T_g)\Delta_{X/Y} \colon X \to X \times_Y X
$$
et par suite, dire que $X$ est formellement principal homog\`ene
sous $G_Y$ signifie aussi que $X \times_Y X$ est isomorphe \`a la
somme directe des transform\'ees de la diagonale par les
\'el\'ements $(e,g)$ de~$G \times G$ (op\'erant sur~$X \times_Y
X$ de fa\c con \'evidente) o\`u $e$ d\'esigne l'\'el\'ement
unit\'e de~$G$. Si on ne veut pas distinguer la gauche et la droite
et donner une formule qui reste applicable au produit de plus de deux
facteurs identiques \`a $X$, on peut formuler la condition en disant
que le morphisme canonique
$$
X \times_G (G \times G) \to X \times_Y X
$$
obtenu en attachant au couple $(g,g')$ le morphisme
$$
(T_g,T_{g'})=(T_g \times_Y T_{g'})\Delta_{X/Y} \colon X \to X \times_Y
X
$$
et en faisant op\'erer $G$ \`a gauche sur~$G \times G$ par
l'homomorphisme diagonal;
$$
s(g,g')=(sg,sg'),
$$
est un \emph{isomorphisme}.

La
notion d'\emph{espace principal homogène}
\index{principal homog\`ene (pr\'esch\'ema, fibr\'e)|hyperpage}est d\'eduite de celle de l'espace formellement principal
homog\`ene en ajoutant un axiome suppl\'ementaire, assurant que le
\og quotient\fg de~$X$ par $G_Y$ existe et est pr\'ecis\'ement
l'objet unit\'e \`a droite de la cat\'egorie, ici $Y$. Cet
axiome peut varier suivant le contexte, et s'explicite souvent le plus
commod\'ement (dans le yoga de la \og descente\fg) en exigeant que
l'objet \`a op\'erateurs devienne \og trivial\fg \ie isomorphe au
produit $X \times_Y G_Y$ (en l'occurrence $X \times G$) par un
changement de base convenable, de type pr\'ecis\'e (de telle fa\c con, en pratique, \`a permettre la technique de descente;
\cf Grothendieck, Technique de descente et th\'eor\`emes
d'existence en G\'eom\'etrie Alg\'ebrique, S\'em. Bourbaki
\No 190, pages 26 \`a 28)\footnote{\Cf Exp.\, \ref{VIII} pour la
th\'eorie de la descente plate.}. Dans cet ordre d'id\'ees,
signalons ici la caract\'erisation des fibr\'es principaux
homog\`enes de groupe $G$ (au sens de \loccit):

\begin{proposition}
\label{V.2.6}
Soient $Y$ un pr\'esch\'ema localement noeth\'erien, $X$ un
$Y$-pr\'esch\'ema \`a groupe fini $G$ d'op\'erateurs
op\'erant \`a droite. Les conditions suivantes sont
\'equivalentes:
\begin{enumerate}
\item[(i)] $X$ est fini sur~$Y$, $Y=X/G$, les groupes d'inertie des
points de~$X$ sont r\'eduits \`a l'unit\'e.
\item[(ii)] Il existe un changement de base fid\`element plat et
quasi-compact $Y_1 \to Y$ tel que $X_1=X \times_Y Y_1$ soit un
$Y_1$-pr\'esch\'ema \`a op\'erateurs trivial, \ie isomorphe
\`a $Y_1 \times G$.
\item[(ii~bis)] Comme (ii), mais $Y_1 \to Y$ \'etant fini,
\'etale, surjectif.
\item[(iii)] $X$ est formellement principal homog\`ene sous $G_Y$,
et fid\`element plat et quasi-compact sur~$Y$.
\end{enumerate}
\end{proposition}
\textit{Démonstration} (i)$\To$(ii~bis) On prendra
$Y_1=X$, notant que $X \to Y$ est bien fini, \'etale par~\ref{V.2.3}
et surjectif. Montrons que $X_1$ est alors trivial sur~$Y_1$, ce que
r\'esultera du

\begin{corollaire}
\label{V.2.7} Si \textup{(i)} est v\'erifi\'e et si $X$ admet une
section sur~$Y$, alors $X$ est un espace \`a op\'erateurs trivial.
\end{corollaire}

En effet, cette section permet de d\'efinir un $G$-morphisme $X
\times G \to X$, surjectif puisque $G$ est transitif sur les fibres de
$X$, injectif puisque $G$ op\`ere sans inertie; enfin, c'est un
isomorphisme local en vertu de I~\ref{I.5.3} puisque $X$ est \'etale
sur~$Y$. Donc c'est un isomorphisme.

(ii~bis)
implique trivialement (ii), qui implique (i) car les ingr\'edients
de (i) sont \og invariants\fg par extension fid\`element plate
quasi-compacte de la base (\cf S\'eminaire Bourbaki cit\'e plus
haut pour \og fini\fg; pour les groupes d'inertie, on
applique~\ref{V.2.1}, et pour $Y=X/G$, une r\'eciproque
de~\ref{V.1.9} dans le cas d'un changement de base
\emph{fidèlement plat}, que nous avions oubli\'e d'expliciter).

Nous avons prouv\'e (i)$\To$(iii) en passant en prouvant
(i)$\To$(ii~bis). Enfin (iii)$\To$(ii), car la
premi\`ere hypoth\`ese dans (iii) signifie pr\'ecis\'ement que
$X$ devient trivial en faisant le changement de base $Y_1=X$; d'o\`u
(ii) puisque $X$ est fid\`element plat et quasi-compact sur~$Y$.

\begin{definition}
\label{V.2.8} Un $Y$-pr\'esch\'ema $X$ \`a groupe
d'op\'erateurs \`a droite $G$ satisfaisant les conditions
\'equivalentes de~\ref{V.2.6} est appel\'e un \emph{revêtement
principal de~$Y$, de groupe de Galois~$G$}.
\index{groupe de Galois d'un rev\^etement principal|hyperpage}\index{principal (revetement)@principal (rev\^etement)|hyperpage}\index{revetement principal@rev\^etement principal|hyperpage}\end{definition}

\section{Automorphismes et morphismes de rev\^etements \'etales}
\label{V.3}

\begin{proposition}
\label{V.3.1} Soit $X$ \'etale s\'epar\'e de type fini sur~$Y$ localement noeth\'erien, soit $G$ un groupe fini op\'erant sur~$X$ par $Y$-automorphismes. Alors $G$ op\`ere de fa\c con
admissible et le pr\'esch\'ema quotient $X/G$ est \'etale sur~$Y$.
\end{proposition}

On ne suppose pas $X$ fini sur~$Y$, cependant $X$ est quasi-projectif
sur~$Y$ d'o\`u l'existence de~$X/G$ gr\^ace
\`a~\ref{V.1.8}. Prouvons d'abord

\begin{corollaire}
\label{V.3.2} Le morphisme $X \to X/G$ est \'etale.
\end{corollaire}

On peut supposer \'evidemment $G$ transitif sur l'ensemble des
composantes connexes de~$X$, puis par consid\'eration du
stabilisateur d'une composante connexe, que $X$ est m\^eme
connexe. Enfin, on peut supposer que $G$ op\`ere fid\`element.
Mais alors on voit comme dans~\ref{V.2.4} que $G$ op\`ere sans
inertie, donc par~\ref{V.2.3} il s'ensuit que $X \to X/G$ est
\'etale. On conclut gr\^ace au

\begin{lemme}[remords \`a l'expos\'e I]
\label{V.3.3} Soient $X \to X' \to
Y$ des morphismes de type fini, $x$ un point de~$X$, $x'$ et $y$ ses
images. On suppose $Y$ localement noeth\'erien. Si deux des
morphismes envisag\'es sont \'etales aux points marqu\'es, il en est
de m\^eme du troisi\`eme.
\end{lemme}

Il reste seulement \`a regarder le cas o\`u $X \to X'$ et $X \to
Y$ sont \'etales en $x$ et prouver que $X' \to Y$ l'est en $x'$ (ce
qui est le cas dont nous avons besoin
pour~\ref{V.3.1}). Faisant une extension plate convenable de la base
$Y$, on est ramen\'e au cas o\`u l'extension r\'esiduelle
$\kres(x)/\kres(y)$
est triviale. Consid\'erons les homomorphismes
$\cal{O}_y \to \cal{O}_{x'} \to \cal{O}_x$ et les homomorphismes
d\'eduits par passage aux compl\'et\'es, l'hypoth\`ese
signifie que $\hat{\cal{O}_y} \to \hat{\cal{O}_x}$ et
$\hat{\cal{O}_{x'}} \to \hat{\cal{O}_x}$ sont des isomorphismes,
d'o\`u aussit\^ot que $\hat{\cal{O}_y} \to \hat{\cal{O}_{x'}}$ en
est un, ce qui prouve le lemme.

\begin{corollaire}
\label{V.3.4} Si $X$ est fini et \'etale sur~$Y$, alors $X/G$
est fini et \'etale sur~$Y$.
\end{corollaire}

\begin{proposition}
\label{V.3.5} Soient $X$, $X'$ deux rev\^etements \'etales de
$Y$. Alors tout $Y$-morphisme $f \colon X \to X'$ se factorise en le
produit d'un morphisme \'etale surjectif $X \to X''$ et de
l'immersion canonique $X'' \to X'$ d'une partie $X''$ de~$X'$ \`a la
fois ouverte et ferm\'ee.
\end{proposition}

On sait (I~\ref{I.4.8}) que $f$ est \'etale, donc un morphisme
ouvert, d'autre part $X$ \'etant fini sur~$Y$, $f$ est ferm\'e,
donc $f(X)=X''$ est une partie \`a la fois ouverte et ferm\'ee de~$X'$. On a fini (N.B. il suffisait que $X'$, au lieu d'un
rev\^etement \'etale, soit non ramifi\'e sur~$Y$).

\begin{corollaire}
\label{V.3.6} Avec les notations pr\'ec\'edentes, $X \to X'$
est un \'epimorphisme strict dans la cat\'egorie des
pr\'esch\'emas, et $X' \to X''$ est un monomorphisme (et m\^eme
un monomorphisme strict) dans la cat\'egorie des pr\'esch\'emas.
\end{corollaire}

La premi\`ere assertion signifie par d\'efinition que la suite de
morphismes
$$
\xymatrix@C=.7cm{{X \times_{X''} X}\ar@<2pt>[r]^-{\pr_1}\ar@<-2pt>[r]_-{\pr_2}&{X}\ar[r]&{X''}}
$$
est exacte, et cela r\'esulte du fait que $X \to X''$ est fini et
fid\`element plat, comme on voit facilement (\cf Grothendieck, \loccit). L'assertion duale pour $X'' \to X'$ est encore plus triviale.

Le corollaire~\ref{V.3.6} nous sera utile pour la th\'eorie du
groupe fondamental au \No suivant; il est possible (pour ceux
qui n'aiment pas la notion d'\'epimorphisme strict) de remplacer le
corollaire~\ref{V.3.6} par telle variante que le lecteur arrangera
\`a son go\^ut personnel. Profitons seulement de l'occasion pour
signaler qu'une factorisation
$f=f'f''$, avec $f''$ un \'epimorphisme strict et $f'$ un
monomorphisme, est n\'ecessairement unique \`a isomorphisme unique
pr\`es (dans toute cat\'egorie); cependant, il peut exister en
m\^eme temps une factorisation $f=f_1f_2$ ayant les
propri\'et\'es duales: $f_2$ est un \'epimorphisme, $f_1$ un
monomorphisme strict, (\'egalement unique \`a isomorphisme unique
pr\`es), qui ne soit pas isomorphe \`a la pr\'ec\'edente: il
suffit de prendre par exemple la cat\'egorie des espaces vectoriels
topologiques (s\'epar\'es, si on y tient), et pour $u \colon X \to
X'$ un morphisme tel que $u(X)$ ne soit pas ferm\'e.

\begin{proposition}
\label{V.3.7} Soient $Y$ un pr\'esch\'ema \emph{connexe} localement
noeth\'erien, $y$ un point de~$Y$, $\mathit{\Omega}$ une extension
alg\'ebriquement close
de~$\kres(y)$. Pour tout $X$ sur~$Y$, on
d\'esigne par $X(\Omega)$ l'ensemble des points de~$X$
\`a valeurs dans $\Omega$. Soient $X$, $X'$ des rev\^etements
\'etales de~$Y$ et $u \colon X \to X'$ un $Y$-morphisme tel que
l'application correspondante $X(\Omega) \to
X'(\Omega)$ soit un isomorphisme. Alors $u$ est un
isomorphisme.
\end{proposition}

On est imm\'ediatement ramen\'e au cas o\`u $X'$ est
connexe. Comme $X \to X'$ est fini et \'etale, on sait que le nombre
g\'eom\'etrique de points dans une fibre de~$X \to X'$ est
constant, et \'egal \`a 1 si et seulement si le morphisme
consid\'er\'e est un isomorphisme. Or ce nombre est aussi le
nombre d'\'el\'ements dans une fibre de~$X(\Omega) \to
X'(\Omega)$, d'o\`u la conclusion.
\enlargethispage{.5cm}
\section{Conditions axiomatiques d'une th\'eorie de Galois}
\label{V.4}
Soit $\cal{C}$ une cat\'egorie, $F$ un foncteur covariant de
$\cal{C}$ dans la cat\'egorie des ensembles finis. Supposons les
conditions suivantes satisfaites:

\begin{enumerateb}

\item[(G 1)] $\cal{C}$ a un objet final\footnote{Rappelons qu'un objet
$e$ de~$\cal{C}$ est appel\'e \emph{objet final} si pour tout $X$ dans
$\cal{C}$, $\Hom(X,e)$ a exactement un \'el\'ement. On d\'efinit
de fa\c con duale un \emph{objet initial} de~$\cal{C}$.} et le produit
fibr\'e de deux objets au-dessus d'un troisi\`eme dans $\cal{C}$
existe (cet axiome peut aussi s'\'enoncer en disant que dans
$\cal{C}$ les limites projectives finies existent).

\item[(G 2)] Les sommes finies dans $\cal{C}$ existent (donc aussi un
objet initial $\emptyset_{\cal{C}}$, jouant le r\^ole de l'ensemble
vide), ainsi que le quotient d'un objet de~$\cal{C}$ par un groupe
fini d'automorphismes.

\item[(G 3)] Soit $u \colon X \to Y$ un morphisme dans $\cal{C}$,
alors $u$ se factorise en un produit $X\lto{u'} Y'\lto{u''} Y$, avec $u'$ un \'epimorphisme \emph{strict} et
$u''$ un monomorphisme, qui est un isomorphisme sur un sommande direct
de~$Y$.

\item[(G 4)]
Le foncteur $F$ est exact \`a gauche (\ie transforme unit\'e
\`a droite en unit\'e \`a droite, et commute aux produits
fibr\'es).

\item[(G 5)] $F$ commute aux sommes directes finies, transforme
\'epimorphismes stricts en \'epi\-mor\-phismes, et commute au
passage au quotient par un groupe fini d'automorphismes.

\item[(G 6)] Soit $u \colon X \to Y$ un morphisme dans $\cal{C}$ tel
que $F(u)$ soit un isomorphisme, alors $u$ est un isomorphisme.
\end{enumerateb}

Notre objet est construire un groupe topologique~$\pi$, limite
projective de groupes finis, et une \'equivalence de la
cat\'egorie $\cal{C}$ avec la cat\'egorie $\cal{C}(\pi)$
\label{indnot:eb}\emph{des ensembles finis où} $\pi$ \emph{opère
continûment} (\ie de telle fa\c con que le stabilisateur d'un
point soit un sous-groupe ouvert, ou encore qu'il existe un groupe
quotient discret qui op\`ere d\'ej\`a sur l'ensemble
envisag\'e), l'\'equivalence construite transformant le foncteur
donn\'e $F$ en le foncteur d'inclusion \'evident de~$\cal{C}(\pi)$
dans la cat\'egorie des ensembles finis. On notera tout de suite que
la cat\'egorie $\cal{C}(\pi)$ construite \`a l'aide d'un groupe
topologique~$\pi$, et le foncteur d'inclusion pr\'ec\'edent,
satisfont bien aux conditions (G 1) \`a (G 6).

Nous proc\'edons en plusieurs \'etapes.
\begin{enumerateb}
\item[a)] \emph{Soit} $u \colon X \to Y$ \emph{dans}
$\cal{C}$. \emph{Pour que} $u$ \emph{soit un monomorphisme, il faut et
il suffit} $F(u)$ \emph{le soit}. (Utilise (G 1), (G 4), (G 6)).

En effet, dire que $u$ est un monomorphisme signifie que la projection
$\pr_1 \colon X \times_Y X \to X$ est un isomorphisme.

\item[b)] \emph{Tout objet} $X$ \emph{de} $\cal{C}$ \emph{est
artinien}.

En effet, si $X' \to X'' \to X$ sont des monomorphismes tels que
$F(X')$ et $F(X'')$ aient m\^eme image dans $F(X)$, alors par a)
$F(X') \to F(X'')$ est un isomorphisme, donc $X' \to X''$ est un
isomorphisme par (G 6).

\item[c)] \emph{Le foncteur} $F$ \emph{est strictement
pro-représentable}. (\cf Grothendieck, Technique de descente et
th\'eor\`emes d'existence en G\'eom\'etrie Alg\'ebrique, II,
S\'eminaire Bourbaki 195, f\'evrier 1960).

En effet, d'apr\`es \loccit prop.\, \ref{V.3.1}, cela r\'esulte de
b) et (G 4). On peut donc trouver un syst\`eme projectif sur~$I$
ordonn\'e filtrant:
$$
P=(P_i)_{i \in I}
$$
dans $\cal{C}$, consid\'er\'e comme pro-objet de~$\cal{C}$, et un
isomorphisme fonctoriel
\begin{equation*}
\label{eq:V.4.*} \tag{$*$} F(X)=\Hom_{\Pro(\cal{C})}(P,X)=\varinjlim_{i} \Hom_{\cal{C}}(P_i,X)
\end{equation*}
Cet isomorphisme est r\'ealis\'e par un \'el\'ement
$$
\varphi \in \varprojlim_{i} F(P_i)=F(P)
$$
On peut supposer de plus que les homomorphismes de transition
$\varphi_{ji} \colon P_i \to P_j$ ($i \geq j$) sont des
\emph{épimorphismes}, et que \emph{tout épimorphisme} $P_i \to
P'$ soit \'equivalent \`a un \'epimorphisme $P_i \to P_j$ pour
$j \leq i$ convenable (ce qui d\'etermine le syst\`eme projectif~$P$ de fa\c con essentiellement unique).

Un objet $X \in \cal{C}$ est dit \emph{connexe}
\index{connexe (objet d'une cat\'egorie)|hyperpage}s'il n'est pas isomorphe \`a la somme de deux objets de~$\cal{C}$
non isomorphes \`a l'objet initial $\emptyset_{\cal{C}}$.

\item[d)] \emph{Les} $P_i$ \emph{sont connexes et non isomorphes
à} $\emptyset_{\cal{C}}$.

Si $X$ est une unit\'e \`a gauche, on a $F(X)=\emptyset$ par (G 5)
appliqu\'e \`a la somme directe d'une famille vide, et
r\'eciproquement par (G 6). Donc si $X'$ est un objet de~$\cal{C}$
qui n'est pas unit\'e \`a gauche, \ie tel que $F(X') \neq
\emptyset$, il n'existe aucun morphisme de~$X'$ dans $X$. Donc si un
$P_i$ est unit\'e \`a gauche, alors $i$ est un plus grand
\'el\'ement de l'ensemble d'indices ordonn\'e filtrant $I$, et
la formule \eqref{eq:V.4.*} signifierait $F(X)=\Hom(P_i,X)={}$ ensemble
r\'eduit \`a un \'el\'ement pour tout $X$, ce qui est absurde
puisque $F(\emptyset_{\cal{C}})=\emptyset$. Donc les~$P_i$ sont non
isomorphes \`a $\emptyset_{\cal{C}}$.

Supposons qu'on ait $P_i=A \amalg B$, d'o\`u par (G 5)
$F(P_i)=F(A) \amalg F(B)$, en particulier l'\'el\'ement $a_i$
de~$F(P_i)$, correspondant par \eqref{eq:V.4.*} \`a
l'homomorphisme identique $P_i \to P_i$, est dans $F(A) \amalg
F(B)$, par exemple dans $F(A)$. Cela signifie qu'il existe un $j
\geq i$ tel que $\varphi_{ij} \colon P_j \to P_i$ se factorise en
$P_j \to A \to P_i=A \amalg B$, o\`u la deuxi\`eme fl\`eche
est le morphisme canonique. Donc $F(P_j) \to F(P_i)$ se factorise
en $F(P_j) \to F(A) \to F(P_i)=F(A) \amalg F(B)$, et comme $F(P_j)
\to F(P_i)$ est surjectif par (G 5), il s'ensuit que
$F(B)=\emptyset$, donc $B$ est isomorphe \`a $\emptyset_{\cal{C}}$.

\item[e)]
\emph{Tout morphisme} $u \colon X \to Y$ \emph{dans} $\cal{C}$,
\emph{avec} $X$ \emph{non isomorphe à} $\emptyset_{\cal{C}}$
\emph{et} $Y$ \emph{connexe, est un épimorphisme strict. Tout
endomorphisme d'un objet connexe est un automorphisme}.

Consid\'erons la factorisation (G 3) de~$u$, comme $X \neq
\emptyset_{\cal{C}}$ il r\'esulte de (G 6) que $F(X) \neq \emptyset$ donc
$F(Y') \neq \emptyset$ donc $Y' \neq \emptyset_{\cal{C}}$, donc $Y$
\'etant connexe, $Y'$ s'identifie \`a~$Y$, donc $u$ est un
\'epimorphisme strict. Supposons que $u$ soit un endomorphisme de
l'objet connexe $X$, prouvons que c'est un automorphisme. En effet, on
peut supposer~$X$ non isomorphe \`a $\emptyset_{\cal{C}}$, donc $u$ est
un \'epimorphisme strict par ce qui pr\'ec\`ede, donc $F(u)$ est
un \'epimorphisme par (G 5), et comme $F(X)$ est un ensemble fini,
$F(u)$ est bijectif. Donc $u$ est un automorphisme par (G 6).

En particulier, \emph{tout endomorphisme d'un} $P_i$ \emph{est un
automorphisme}.

\item[f)] \emph{Les conditions suivantes sur un} $P_i$ \emph{sont
équivalentes}: (i) L'application injective naturelle
$\Hom(P_i,P_i) \to \Hom(P,P_i) \simeq F(P_i)$ est aussi surjective,
\ie pour tout $u \colon P \to P_i$ il existe un $v \colon P_i \to
P_i$ tel que $u=v\varphi_i$ (o\`u $\varphi_i$ est l'homomorphisme
canonique $P \to P_i$). (ii) Le groupe $\Aut(P_i)$ op\`ere de fa\c con transitive sur~$F(P_i)$. (iii)~Le groupe $\Aut(P_i)$ op\`ere de
fa\c con simplement transitive sur~$F(P_i)$.

En effet, identifiant $\Hom(P,P_i)$ \`a $F(P_i)$, l'application
envisag\'ee dans (i) n'est autre que $v \mto
F(v)(\varphi_i)$. L'\'equivalence des trois conditions provient
alors du fait que $\Hom(P_i,P_i)=\Aut(P_i)$ et que l'application
pr\'ec\'edente est d\'ej\`a injective.

Un $P_i$ satisfaisant les conditions \'equivalentes
(i), (ii), (iii)
de~f) est appel\'e \emph{galoisien}.
\index{galoisien (objet)|hyperpage}
\item[g)] \emph{Pour tout} $X$ \emph{dans} $\cal{C}$, \emph{il existe
un} $P_i$ galoisien \emph{tel que tout} $u \in \Hom(P,X)$ \emph{se
factorise en} $P \lto{\varphi_i} P_i \longrightarrow X$.

Soit $J=\Hom(P,X)=F(X)$, c'est un ensemble fini, donc il existe un
$P_j$ tel que tout $u \colon P \to X$ se factorise en $P\lto{j}P_j\longrightarrow X$, ou encore tel que le morphisme naturel
$$
P \to X^J \qquad (J=\Hom(P,X))
$$
se factorise en
$$
P \lto{\varphi_j} P_j \to X^J
$$
En
vertu de (G~3), le morphisme $P_j \to X^J$ se factorise en le produit
d'un monomorphisme par un \'epimorphisme strict, que l'on peut
prendre de la forme $\varphi_{ij} \colon P_j \to P_i$. On est donc
ramen\'e \`a prouver que $P_i$ est galoisien. Soit $k$ un indice
$\geq j$ tel que tout morphisme $P \to P_i$ se factorise par
$P \lto{\varphi_k} P_k \longrightarrow P_i$. Notons que le
morphisme naturel $P_k \to X^J$ se factorise encore en le compos\'e
$$
P_k \lto{\varphi_{ik}} P_i \lto{U} X^J
$$
o\`u la premi\`ere fl\`eche est un \'epimorphisme strict
par~e), et la deuxi\`eme un monomorphisme. On veut prouver que pour
un morphisme donn\'e $\psi \colon P_k \to P_i$, il existe un
endomorphisme $v$ de~$P_i$ tel que $\psi = v \varphi_{ik}$. Mais pour
tout $u \in \Hom(P_i,X)$, consid\'erons $u \psi \in \Hom(P_k,X)$,
il est donc de la forme $u'\varphi_{ik}$ avec un $u' \in \Hom(P_i,X)$
bien d\'etermin\'e. L'application $u \mto u'$ de~$J$ dans $J$
ainsi d\'efinie par $\psi$ est d'ailleurs injective car $\psi$ est
un \'epimorphisme en vertu de e); elle est donc bijective puisque
l'ensemble~$J$ est fini. L'application bijective $u \mto u'$ de~$J$
dans $J$ d\'efinit donc un isomorphisme $\alpha \colon X^J \isomto
X^J$ rendant commutatif le diagramme
$$
\xymatrix{{P_k}\ar[r]^-{\varphi_{ik}}\ar@{=}[d]&{P_i}\ar[r]^-U&{X^J}\ar[d]^-{\alpha}_-{\simeq}\\{P_k}\ar[r]^{\psi}&{P_i}\ar[r]^-U&{X^J}}
$$
D'apr\`es les propri\'et\'es d'unicit\'e de la factorisation
d'un morphisme en produit d'un monomorphisme par un \'epimorphisme
strict, il s'ensuit (puisque $\psi$ lui aussi est un \'epimorphisme
strict par e)) que l'on peut trouver un morphisme $v \colon P_i \to
P_i$ qui laisse le diagramme commutatif, cqfd.

On en conclut en particulier que \emph{les} $P_i$ \emph{galoisiens
forment un système cofinal dans le système des} $(P_j)$. On
aura donc, puisque pour un objet galoisien $P_i$ on a
$$
\Hom(P,P_i)=\Hom(P_i,P_i)=\Aut(P_i),
$$
par passage \`a la limite:
$$
\Hom(P,P) = \varprojlim_{i} \Hom(P,P_i) = \varprojlim_{i} \Hom(P_i,P_i) =
\varprojlim_{i} \Aut(P_i)
$$
o\`u
\enlargethispage{.5cm}
la limite projective est prise sur les $P_i$ galoisiens. D'ailleurs,
moyennant l'identification $\Hom(P,P_i)=F(P_i)$, et compte tenu que
$F$ transforme \'epimorphismes en \'epimorphismes, on voit que les
homomorphismes de transition dans le syst\`eme projectif
pr\'ec\'edent sont surjectifs. On conclut de tout ceci:

\item[h)] \emph{On a}
$$
\Hom(P,P) = \Aut(P) = \varprojlim_{i} F(P_i) = \varprojlim_{i} \Aut(P_i),
$$
\emph{où la limite projective est prise sur les} $P_i$
\emph{galoisiens}.

En particulier, $\Aut(P)$ appara\^it comme limite projective d'un
syst\`eme projectif de groupes finis (les homomorphismes de
transition \'etant surjectifs), on le munira de la topologie limite
projective des topologies discr\`etes. \emph{On désignera
par}~$\pi$
\label{indnot:ec}\emph{et on appellera groupe fondamental}
\index{groupe fondamental d'une cat\'egorie galoisienne|hyperpage}(de~$\cal{C}$ muni de~$F$) \emph{le groupe opposé
à}~$\Aut(P)$. Ce groupe op\`ere donc \emph{à droite}
sur~$P$, c'est la limite projective de groupes finis $\pi_i$
op\'erant \`a droite sur les $P_i$ galoisiens, $\pi_i$ \'etant
le groupe oppos\'e \`a $\Aut(P_i)$.

Compte tenu de l'isomorphisme fonctoriel
$$
F(X)=\Hom(P,X)
$$
et de la d\'efinition de~$\pi$, on voit donc que $\pi$ op\`ere
\emph{à gauche} sur~$F(X)$, et d'ailleurs de fa\c con continue
d'apr\`es g) (car avec les notations de g), c'est en fait $\pi_i$
qui op\`ere sur~$F(X)$). Il est trivial que pour tout morphisme $u
\colon X \to Y$ dans $\cal{C}$, le morphisme $F(u) \colon F(X) \to
F(Y)$ est compatible avec les op\'erations de~$\pi$. \emph{On peut
donc considérer par la suite} $F$ \emph{comme un foncteur
covariant}
$$
F \colon \cal{C} \to \cal{C}(\pi)
$$
\emph{où} $\cal{C}(\pi)$ \emph{est la catégorie des ensembles
finis où} $\pi$ \emph{opère à gauche continûment}.

Nous allons maintenant d\'efinir un foncteur en sens inverse:
$$
G \colon \cal{C} \from \cal{C}(\pi)
$$
par la formule
$$
G(E)=P \times_{\pi} E,
$$
o\`u $P \times_{\pi} E$ est d\'efini comme solution du
probl\`eme universel r\'esum\'e par
$$
\Hom_{\cal{C}}(P \times_{\pi} E,X) \isomto \Hom_{\pi}(E,\Hom(P,X))
$$
(o\`u
dans le deuxi\`eme membre $\Hom(P,X)=F(X)$ est consid\'er\'e
comme un ensemble o\`u $\pi$ op\`ere \`a gauche). Il faut prouver
l'existence de l'objet $P \times_{\pi} E$.

\item[i)] \emph{Soit} $Q$ \emph{un objet de} $\cal{C}$ \emph{où un
groupe fini} $G$ \emph{opère à droite, et} $E$ \emph{un
ensemble fini où} $G$ \emph{opère à gauche. Alors} $G
\times_G E$ \emph{existe, et l'application canonique}
$$
F(Q) \times_G E \to F(Q \times_G E)
$$
\emph{est un isomorphisme}.

Comme les sommes directes finies existent dans $\cal{C}$ par (G 2), et
que $F$ y commute par (G 5), on est ramen\'e aussit\^ot au cas
o\`u $G$ op\`ere transitivement sur~$E$, car si les~$E_j$ sont les
trajectoires de~$G$ dans $E$, on aura
$$
Q \times_G E = \underset {j}{\amalg}\ Q \times_G E_j.
$$
Soit alors $a \in E$, soit $H$ son stabilisateur, on voit aussit\^ot
sur la d\'efinition que $Q \times_G E$ s'identifie \`a $Q/H$.
D'o\`u l'existence, gr\^ace \`a (G~2), et la propri\'et\'e
de commutation pour~$F$ gr\^ace \`a (G~5).

\item[j)] \emph{Soit} $E$ \emph{un objet de} $\cal{C}(\pi)$, \emph{et
soit} $P_i$ \emph{galoisien tel que} $\pi_i$ \emph{opère
déjà sur} $E$. \emph{Alors} $P_i \times_{\pi_i} E$
\emph{existe et on a un isomorphisme canonique}
$$
E \isomto F(P_i \times_{\pi_i} E)
$$
\emph{Si} $j \geq i$ \emph{est tel que} $P_j$ \emph{soit galoisien,
alors l'homomorphisme canonique} $P_j \times_{\pi_j} E \to P_i
\times_{\pi_i} E$ \emph{est un isomorphisme}.

La premi\`ere assertion est un cas particulier de i), compte tenu
que $\pi_i$ op\`ere de fa\c con simplement transitive sur~$F(P_i)$
qui est muni d'un point marqu\'e $\varphi_i$, d'o\`u un isomorphisme
$F(P_i) \times_{\pi_i} E \simeq E$. Pour la deuxi\`eme assertion on
utilise par exemple (G 6).

Soit, pout tout $j$, $\cal{C}_j$ la sous-cat\'egorie pleine de
$\cal{C}$ form\'ee des $X$ tels que $\Hom(P_j,X) \to \Hom(P,X)
\simeq F(X)$ soit bijectif. On sait par g) que $\cal{C}$ \emph{est
réunion filtrante des} $\cal{C}_j$. On a donc pour $X \in
\cal{C}_j$:
\begin{align*}
\Hom_{\pi}(E,\Hom(P,X)) & \simeq \Hom_{\pi}(E,\Hom(P_j,X)) \simeq
\Hom_{\pi_j}(E,\Hom(P_j,X))\\ & \simeq \Hom(P_j \times_{\pi_j} E,X)
\end{align*}
et
compte tenu de la derni\`ere assertion dans j) on trouve un
isomorphisme fonctoriel en l'objet $X$ de~$\cal{C}_j$:
$$
\Hom_{\pi}(E,\Hom(P,X)) \simeq \Hom(P_i \times_{\pi_i} E,X)
$$
Comme cela est vrai pour tout $j$ et que ces isomorphismes fonctoriels,
pour $j$ variable, s'induisent mutuellement, on conclut:

\item[k)] \emph{Sous les conditions de} j), \emph{le morphisme
composé des morphismes canoniques}
$$
E \to \Hom(P_i,P_i \times_{\pi_i} E) \to \Hom(P,P_i \times_{\pi_i} E)
$$
\emph{fait de} $P_i \times_{\pi_i} E$ \emph{une solution du
problème universel définissant} $P \times_{\pi} E$,
\ie \emph{ce dernier existe et on a un isomorphisme}
$$
P \times_{\pi} E \isomto P_i \times_{\pi_i} E
$$

Cela ach\`eve la construction du foncteur $G(E)$. On a d'autre part
un homomorphisme fonctoriel
$$
\alpha \colon \id_{\cal{C}(\pi)} \to FG
$$
\ie un homomorphisme fonctoriel en l'objet $E$ de~$\cal{C}(\pi)$:
$$
\alpha(E) \colon E \to FG(E) = F(P \times_{\pi} E)
$$
savoir le compos\'e des morphismes canoniques
$$
E \to F(P) \times_{\pi} E \to F(P \times_{\pi} E)
$$
(o\`u le premier provient du point marqu\'e $\varphi \in F(P)$).
Conjuguant j) et k), on trouve:

\item[l)] \emph{L'homomorphisme} $\alpha$ \emph{est un isomorphisme}

On d\'efinit de m\^eme un homomorphisme fonctoriel
$$
\beta \colon GF \to \id_{\cal{C}}
$$
\ie un homomorphisme fonctoriel en l'objet $X$ de~$\cal{C}$:
$$
\beta(X) \colon P \times_{\pi} F(X) \to X
$$
comme
associ\'e au $\pi$-homomorphisme
$$
F(X) \to \Hom(P,X)
$$
inverse de l'isomorphisme canonique $\Hom(P,X) \isomto F(X)$.

\item[m)] \emph{Les composés}
\begin{gather*}
F(X) \lto{\alpha(F(X))} FGF(X) \lto{F(\beta(X))} F(X)\\
G(E) \lto{G(\alpha(E))} GFG(E) \lto{\beta(G(E))} G(E)
\end{gather*}
\emph{sont les isomorphismes identiques}.

\^Ane qui trotte.

Compte tenu de l) il s'ensuit:

\item[n)] \emph{L'homomorphisme} $\beta$ \emph{est un isomorphisme}

Nous avons ainsi obtenu le r\'esultat promis:
\end{enumerateb}

\begin{theoreme}
\label{V.4.1}
Soit $\cal{C}$ une cat\'egorie satisfaisant les conditions \emph{(G
1)}, \emph{(G 2)}, \emph{(G~3)} du d\'ebut du num\'ero, et $F$ un
foncteur covariant de~$\cal{C}$ dans la cat\'egorie des ensembles
finis, satisfaisant les conditions \emph{(G 4)}, \emph{(G 5)} et
\emph{(G 6)}. Alors les constructions canoniques pr\'ec\'edentes
d\'efinissent des \'equivalences de cat\'egories $F\colon
\cal{C}\to \cal{C}(\pi)$ et $G\colon \cal{C}(\pi)\to \cal{C}$
quasi-inverses l'une de l'autre. De fa\c con pr\'ecise, il existe
un pro-objet~$P$ de~$\cal{C}$, et un isomorphisme fonctoriel $F(X)
\isomfrom \Hom(P, X)$, $\pi$ est le groupe oppos\'e au groupe des
automorphismes de~$P$, topologis\'e de fa\c con convenable, de
fa\c con que~$\pi$ op\`ere de fa\c con continue sur les
ensembles $\Hom(P, X)\simeq F(X)$. Enfin $G$ est donn\'e par
$G(E)\simeq P\times_{\pi} E$.
\end{theoreme}

\begin{remarques}
\label{V.4.2}
L'\'enonc\'e des conditions (G 1) \`a (G 6) devient plus simple
et plus sympathique si on remplace (G 2) et (G 5) respectivement par:
\begin{enumerateb}
\item[(G' 2)] Les limites inductives finies dans
$\cal{C}$
existent.
\item[(G' 5)] Le foncteur $F$ est exact \`a droite (\ie commute aux
limites inductives finies).
\end{enumerateb}
Ces conditions sont en apparence plus fortes que (G~2) et (G~5), mais
il r\'esulte aussit\^ot du th\'eor\`eme de
structure~\ref{V.4.1} qu'elles sont entra\^in\'ees par (G~1)
\`a (G~6). On notera cependant que dans les cas qui nous
int\'eresseront, la v\'erification de (G~2) et (G~5) semble
effectivement plus simple que celle de (G'~2) et (G'~5). J'ignore si
dans la condition (G~3), le fait que $u''$ soit un isomorphisme sur un
sommande direct de~$Y$ ne pourrait \^etre omis.
\end{remarques}

\section{Cat\'egories galoisiennes}
\label{V.5}

\begin{definition}
\label{V.5.1}
On appelle cat\'egorie galoisienne
\index{galoisienne (cat\'egorie)|hyperpage}une cat\'egorie $\cal{C}$ \'equivalente \`a une cat\'egorie
$\cal{C}(\pi)$, o\`u $\pi$ est un groupe compact, limite projective
de groupes finis (\ie totalement disconnexe).

Pour la d\'efinition de~$\cal{C}(\pi)$, \cf d\'ebut du \No \ref{V.4}. En
vertu du th.\, \ref{V.4.1}, $\cal{C}$ est galoisienne si et seulement si
elle satisfait les conditions (G 1) \`a (G 3), et s'il existe un
foncteur $F$ de
$\cal{C}$
dans la cat\'egorie des ensembles finis satisfaisant les conditions
(G~4) \`a (G~6) (\ie qui est \emph{exact} et \emph{conservatif},
dans une terminologie g\'en\'erale). Un tel foncteur sera
appel\'e \emph{foncteur fondamental}
\index{foncteur fondamental|hyperpage}de la cat\'egorie galoisienne $\cal{C}$\footnote{Il semble pr\'ef\'erable d'adopter le terme plus parlant
de \og foncteur fibre\fg.};
\index{foncteur fibre|hyperpage}il est pro-repr\'esentable par un pro-objet que nous noterons~$P_F$;
un pro-objet $P$ tel que le foncteur $F$ associ\'e soit fondamental
est appel\'e \emph{pro-objet fondamental}.
\index{pro-objet fondamental|hyperpage}De cette fa\c con, la cat\'egorie des foncteurs fondamentaux sur~$\cal{C}$ est anti-\'equivalente \`a la cat\'egorie des
pro-objets fondamentaux; si $F$ et $P$ se correspondent, le groupe
$\Aut F$ est donc isomorphe \`a l'oppos\'e du groupe $\Aut P$,
donc le groupe not\'e $\pi$ dans le num\'ero pr\'ec\'edent
n'est autre que $\Aut P$. Rappelons qu'au num\'ero pr\'ec\'edent
nous avons construit, \`a partir d'un foncteur fondamental
\emph{donné} $F$, une \'equivalence de~$\cal{C}$ avec
$\cal{C}(\pi)$ (o\`u $\pi=\Aut(F)$) qui transforme le foncteur $F$
en le foncteur canonique de~$\cal{C}(\pi)_x$ dans la cat\'egorie des
ensembles finis. Dans ce cas type $\cal{C}=\cal{C}(\pi)$, $F=$
foncteur canonique, le pro-objet fondamental associ\'e \`a $F$
n'est autre que le syst\`eme projectif des quotients discrets
$\pi_i$ de~$\pi$.
\end{definition}

Il peut \^etre utile d'expliciter la cat\'egorie des pro-objets de
$\cal{C}(\pi)$. On trouve:

\begin{proposition}
\label{V.5.2}
La cat\'egorie $\Prodash\cal{C}(\pi)$
\label{indnot:ed}est canoniquement \'equivalente \`a la cat\'egorie
$\cal{C}'(\pi)$ des espaces, \`a groupe topologique $\pi$
d'op\'erateurs, qui sont compacts et totalement disconnexes.
\end{proposition}

Comme cette derni\`ere contient $\cal{C}(\pi)$ comme
sous-cat\'egorie pleine (correspondant aux espaces \`a
op\'erateurs compacts discrets) et que les limites projectives y
existent, on a en tous cas un foncteur canonique $g\colon
\Prodash\cal{C}(\pi) \to \cal{C}'(\pi)$, au syst\`eme projectif
$Q=(Q_i)$ correspondant
l'objet $X=\varprojlim_{i} Q_i$ de~$\cal{C}'(\pi)$. Pour d\'efinir
un foncteur en sens inverse, il suffit de d\'efinir un foncteur
contravariant de~$\cal{C}'(\pi)$ dans la cat\'egorie des foncteurs
$\cal{C}\to \Ens$ exacts \`a gauche, et on prendra pour tout
$X\in\cal{C}'(\pi)$ le foncteur $h(X)(E)= \Hom(X, E)$
(le $\Hom$ \'etant pris dans $\cal{C}'(\pi)$). Il est imm\'ediat
par d\'efinition que les foncteurs $h$ et $g$ sont adjoints l'un de
l'autre, et que $hg$ est canoniquement isomorphe au foncteur identique
de~$\Prodash\cal{C}(\pi)$. Il reste \`a prouver (pour \'etablir
que $g$ et $h$ sont quasi-inverses l'un de l'autre) que tout objet de
$\cal{C}'(\pi)$ est isomorphe \`a un objet de la forme $g(Q)$,
o\`u $Q\in\Prodash{\cal C}(\pi)$, en d'autres termes: \emph{tout
espace $X$ à groupe topologique $\pi$ d'opérateurs, qui est
compact et totalement disconnexe, est isomorphe à une limite
projective d'espaces à opérateurs finis discrets}. Comme $X$
est limite projective de ses quotients finis discrets (en tant
qu'espace topologique sans op\'erateurs), on est ramen\'e \`a
montrer que dans l'ensemble de ces quotients, il y a un syst\`eme
cofinal qui est invariant par $\pi$. Il suffit pour cela de montrer
que pour un tel quotient $X'$, l'ensemble des transform\'es de ce
quotient par les op\'erations de~$\pi$ est fini, (on prendra alors
le sup desdits transform\'es, qui sera un quotient invariant
majorant $X'$). Ou encore, qu'il y a un sous-groupe invariant ouvert
$\pi'$ de~$\pi$ tel que les \'el\'ements de ce sous-groupe
invarient $X'$. Or $X'$ correspond \`a une partition finie de~$X$ en
ensembles ouverts $X_i$. Par raison de continuit\'e et de
compacit\'e de~$\pi$, il existe un voisinage $V$ de
l'\'el\'ement neutre de~$\pi$ tel que $s\in V$ implique
$s\cdot X_i\subset X_i$ pour tout~$i$, donc $s$ invarie~$X'$. Or on sait
que les sous-groupes invariants ouverts de~$\pi$ forment un
syst\`eme fondamental de voisinages de l'\'el\'ement neutre, ce
qui ach\`eve la d\'emonstration.

Remarquons qu'on voit encore plus simplement que la cat\'egorie
$\Ind \cal{C}(\pi)$ est canoniquement \'equivalente \`a la
cat\'egorie des ensembles o\`u $\pi$ op\`ere
contin\^ument. Nous n'en aurons pas besoin ici.

\begin{proposition}
\label{V.5.3}
Soient $\cal{C}$ une cat\'egorie galoisienne, $F$ un foncteur
fondamental
sur~$\cal{C}$,
$P=(P_i)$ le pro-objet associ\'e, normalis\'e de la fa\c con
habituelle. Soit
$X\in\cal{C}$;
pour que $X$ soit connexe, il faut et il suffit que $\pi$ op\`ere
transitivement sur~$E=F(X)$.
\end{proposition}

On est ramen\'e au cas type $\cal{C}=\cal{C}(\pi)$, $F=$ foncteur
canonique, o\`u c'est trivial.

\begin{corollaire}
\label{V.5.4}
Conditions \'equivalentes sur~$X$: \textup{(i)} $X$ est connexe et
$\not\simeq\emptyset_{\cal{C}}$ \textup{(ii)} le groupe $\pi$ est
transitif sur~$E=F(X)$, et $F(X)\ne\emptyset$ \textup{(iii)} $X$ est
isomorphe \`a un $P_i$.
\end{corollaire}

L'\'equivalence
de (i) et (iii) r\'esulte aussi d\'ej\`a facilement de
\No \ref{V.4},~e).

\begin{proposition}
\label{V.5.5}
Soit $Q=(Q_i)_{i\in I}$ un pro-objet
de~$\cal{C}$,
normalis\'e de la fa\c con habituelle, et soit $G$ le foncteur
correspondant $G(X)=\Hom(Q, X)$ de~$\cal{C}$
(dans~$\Ens$).
Les conditions suivantes sont \'equivalentes:
\begin{enumerate}
\item[(i)] $G$ commute aux sommes directes finies
\item[(ii)] $G$ commute \`a la somme de deux objets
\item[(iii)] Les $Q_i$ sont connexes et
$\not\simeq\emptyset_\cal{C}$
\item[(iv)] $Q$ est isomorphe \`a $\pi/H$, o\`u $H$ est un
sous-groupe ferm\'e de~$\pi$.
\item[(v)] Le foncteur $G$ est isomorphe au foncteur $E\mto
E^H$ (ensemble des invariants par $H$) d\'efini par un sous-groupe
ferm\'e $H$ de~$\pi$.
\end{enumerate}
\end{proposition}

N.B. dans l'\'enonc\'e de (iv) et (v), on suppose choisi un
foncteur fondamental, permettant d'identifier $\cal{C}$ \`a la
cat\'egorie~$\cal{C}(\pi)$.

\subsubsection*{D\'emonstration} On peut supposer
$\cal{C}=\cal{C}(\pi)$. L'implication (i)$\To$(ii) est
triviale, (ii)$\To$(iii) se prouve comme la
propri\'et\'e d) du \No \ref{V.4}. Prouvons (iii)$\To$(iv). En
effet, on sait que $\varprojlim_i Q_i$ est non vide comme limite
projective d'ensembles finis non vides. Soit $a$ un point de
$\varprojlim_i Q_i$, il d\'efinit un homomorphisme d'espaces \`a
op\'erateurs
$$
\pi\to Q
$$
qui est \emph{surjectif}, car pour tout $i$ le compos\'e $\pi\to
Q\to Q_i$ l'est, puisque $\pi$ est transitif dans $Q_i$ en vertu
de~\ref{V.5.3}. Si donc $H$ est le sous-groupe de stabilisateur de
$a$, on obtient un isomorphisme $\pi/H\isomto Q$. Les implications
(iv)$\To$(v) et (v)$\To$(i) sont \`a nouveau
triviales.

\begin{proposition}
\label{V.5.6}
Soient $\cal{C}$ une cat\'egorie galoisienne, $P$ un pro-objet
fondamental de~$\cal{C}$ et $F$ le foncteur fondamental
associ\'e. Soit $P'=(P'_i)_{i\in I}$ un pro-objet de~$\cal{C}$, mis
sous forme normale, et $F'$ le foncteur associ\'e $F'(X)=\Hom(P',X)$
de~$\cal{C}$ dans~$\Ens$. Conditions \'equivalentes:
\begin{enumerate}
\item[(i)] $P'\simeq P$, ou encore $F'\simeq F$.
\item[(ii)] $P'$ est fondamental, ou encore $F'$ est
fondamental.
\item[(iii)] $F'$ transforme somme de deux objets en somme, et
$X\not\simeq\emptyset_{\cal{C}}$ implique $F(X)\ne\emptyset$.
\item[(iv)]
Les objets de~$\cal{C}$ connexes et $\ne\emptyset_{\cal{C}}$ sont
exactement les objets isomorphes \`a un~$P'_i$.
\end{enumerate}
\end{proposition}

On a trivialement (i)$\To$(iii) et (i)$\To$(ii),
de plus (ii)$\To$(iv) en vertu de~\ref{V.5.4} (appliqu\'e
\`a $P'$ au lieu de~$P$). De plus (iii) ou (iv) implique en vertu
de~\ref{V.5.5} que $P'$ est de la forme $\pi/H$ o\`u $H$ est un
sous-groupe ferm\'e de~$\pi$. Dans le cas (iii), il existe pour tout
sous-groupe invariant ouvert $\pi'$ de~$\pi$ un $\pi$-homomorphisme
$P'=\pi/H\to \pi/\pi'$, d'o\`u $H\subset \pi'$, d'o\`u $H=(0)$ et
par suite (i), cqfd.

\begin{corollaire}
\label{V.5.7}
Soit $\cal{C}$ une cat\'egorie galoisienne. Les pro-objets
fondamentaux sont isomorphes, les foncteurs fondamentaux sont
isomorphes.
\end{corollaire}

En d'autres termes, \emph{la catégorie des foncteurs fondamentaux
est un groupoïde connexe~$\Gamma$},
\label{indnot:ef}qu'on peut appeler le \emph{groupoïde fondamental}
\index{groupoide fondamental@groupo\"ide fondamental|hyperpage}de la cat\'egorie galoisienne~$\cal{C}$. Si $\cal{C}=\cal{C}(\pi)$,
le groupe des automorphismes d'un objet du groupo\"ide fondamental
est isomorphe \`a $\pi$, cet isomorphisme \'etant bien
d\'etermin\'e \`a automorphisme int\'erieur pr\`es. (N.B. on
appelle \emph{groupoïde} une cat\'egorie o\`u tous les
morphismes sont des isomorphismes, groupo\"ide \emph{connexe}
groupo\"ide dont tous les objets sont isomorphes). Les pro-objets
fondamentaux de~$\cal{C}$ forment un groupo\"ide connexe
\'equivalent \`a l'\emph{opposé} du groupo\"ide
fondamental. Si $F, F'$ sont deux foncteurs fondamentaux, associ\'es
\`a des pro-objets fondamentaux $P, P'$, alors $\Hom(F, F')=\Isom(F,
F')$ est parfois not\'e $\pi_{F', F}$ et joue le r\^ole d'un
\og ensemble de \emph{classes de chemins} de~$F$ \`a~$F'$\fg. En
particulier $\pi_{F, F}=\pi_F$ n'est autre que le \emph{groupe
fondamental de~$\cal{C}$} en $F$ construit dans le num\'ero
pr\'ec\'edent. Quant au pro-objet $P$ associ\'e \`a~$F$, il
joue le r\^ole d'un \emph{revêtement universel en~$F$} de
l'objet final $e_{\cal{C}}$ de~$\cal{C}$.

Il peut \^etre commode d'avoir une description de~$\cal{C}$ (\`a
\'equivalence pr\`es) en termes de son groupo\"ide fondamental
$\Gamma$, sans passer par un choix d'un objet particulier $F$ de ce
dernier. Or \`a tout objet $X$ de~$\cal{C}$ est associ\'e le
foncteur $E_X$ sur le groupo\"ide fondamental, d\'efini par
$$
E_X(F)=F(X),
$$
\`a valeurs dans~$\Ens$. (Un tel foncteur est connu en topologie
sous le nom de \og syst\`eme local\fg sur le groupo\"ide)
$F(X)=E_X(F)$ peut \^etre appel\'e la \emph{fibre} de~$X$ en $F$,
et le foncteur $E_X$ le foncteur-fibre associ\'e \`a~$X$. Le
foncteur $E_X$ a la propri\'et\'e
suivante:
\emph{pour tout $F$, $E_X(F)$ est un ensemble fini
où le groupe topologique $\pi_F=\Aut(F)$ opère continûment}.
Pour un
foncteur covariant donn\'e $\xi$ du groupo\"ide fondamental
dans~$\Ens$, la condition pr\'ec\'edente \'equivaut d'ailleurs
\`a la m\^eme condition pour \emph{un} $F$ fix\'e
quelconque. Ceci pos\'e:

\begin{proposition}
\label{V.5.8}
Le foncteur $X\mto E_X$ est une \'equivalence de la cat\'egorie
$\cal{C}$ avec la cat\'egorie des foncteurs covariants du
groupo\"ide fondamental $\Gamma$ de~$\cal{C}$ dans~$\Ens$, qui
satisfont la condition
mise en \'evidence ci-dessus.
\end{proposition}

En effet, soit $F_0$ un objet du groupo\"ide fondamental, et soit
$\pi_0=\Aut(F_0)$, alors le foncteur $\xi\mto \xi(F_0)$ est une
\'equivalence de la deuxi\`eme cat\'egorie envisag\'ee
dans~\ref{V.5.8} avec la cat\'egorie $\cal{C}(\pi_0)$, comme on
constate aussit\^ot. D'autre part, le compos\'e de ce dernier et
du foncteur $X\mto E_X$ est l'\'equivalence naturelle
$\cal{C}\to\cal{C}(\pi_0)$. Il en r\'esulte que le foncteur
$X\mto E_X$ lui-m\^eme est une \'equivalence.

\begin{corollaire}
\label{V.5.9}
La cat\'egorie $\Prodash\cal{C}$
est \'equivalente canoniquement \`a la cat\'egorie des foncteurs
covariants $\xi$ du groupo\"ide fondamental $\Gamma$ dans la
cat\'egorie des espaces topologiques, satisfaisant la condition:
pour tout objet $F$ de~$\Gamma$, $\xi(F)$ est un espace compact
totalement disconnexe \`a groupe topologique $\pi_F$
d'op\'erateurs.
\end{corollaire}

Ici encore, on peut v\'erifier cette condition sur~$\xi$, il suffit
de la v\'erifier pour \emph{un}~$F$. La d\'emonstration est la
m\^eme que pour~\ref{V.5.8}.

\begin{remarque}
\label{V.5.10}
Soit $(F_s)_{s\in S}$ une famille
d'objets
du groupo\"ide fondamental~$\Gamma$. Posons pour $s, s'\in S$:
$$
\Hom(s, s')=\Hom(F_s, F_{s'})
$$
de sorte que $S$ devient lui-m\^eme un groupo\"ide connexe et
l'application $s\mto F_s$ un foncteur pleinement fid\`ele de~$S$
dans $\Gamma$, soit~$f$. Consid\'erant alors le foncteur $X\mto
E_X\circ f$ de~$\cal{C}$ dans la cat\'egorie de foncteurs
$\Hom(S,\Ens)$, on obtient une variante de~\ref{V.5.8}
(et~\ref{V.5.9}) avec $\Gamma$ remplac\'e par~$S$. L'\'enonc\'e
ainsi obtenu se r\'eduit au th\'eor\`eme~\ref{V.4.1} lorsque $S$
est r\'eduit \`a un point, et n'est autre que~\ref{V.5.8}
lui-m\^eme si $S$ est l'ensemble des objets de~$\Gamma$.
\end{remarque}

Nous allons utiliser~\ref{V.5.9} pour d\'efinir un pro-objet
canonique de~$\cal{C}$. Pour ceci, nous consid\'erons le foncteur
de~$\Gamma$ dans la cat\'egorie des espaces topologiques
(et m\^eme des groupes topologiques):
$$
f\colon F\mto \Aut(F)=\pi_F.
$$
Ce foncteur satisfait la condition envisag\'ee dans~\ref{V.5.8},
l'espace \`a op\'erateurs $f(F)$ sous $\pi_F$ n'est autre que
$\pi_F$, consid\'er\'e comme espace \`a op\'erateurs sous
lui-m\^eme par les automorphismes int\'erieurs. Donc le foncteur
$f$ correspond \`a un pro-objet de~$\cal{C}$ d\'etermin\'e \`a
isomorphisme unique pr\`es, qui est m\^eme un pro-groupe de
$\cal{C}$ et qui est appel\'e le \emph{pro-groupe fondamental
de~${\cal C}$},
\index{pro-groupe fondamental d'une cat\'egorie galoisienne|hyperpage}jouant le r\^ole d'un syst\`eme local de groupes
fondamentaux. C'est donc un pro-groupe $\Pi$ de~$\cal{C}$ d\'efini
par la condition qu'on ait un isomorphisme fonctoriel en $F$
$$
F(\Pi)\simeq\pi_F.
$$

Si $X$ est un pro-objet quelconque de~$\cal{C}$, on a un morphisme
canonique
$$
\Pi\times X \to X
$$
qui fait de~$X$ un objet \`a groupe d'op\'erateurs \`a gauche
$G$ dans $\Prodash\cal{C}$. Il suffit pour ceci de noter que pour $F$
variable, on a une application canonique
$$
\Pi(F)\times X(F) \to X(F)
$$
\ie $$
\Aut(F)\times E_X(F) \to E_X(F), \quad \text{ou} \quad \pi_F\times
F(X)\to F(X)
$$
qui est fonctorielle en~$F$. Elle est aussi fonctorielle en $X$, donc
pour tout morphisme $X\to Y$ de pro-objets, le diagramme
$$
\xymatrix{{\Pi \times X}\ar[r]\ar[d]&{X}\ar[d]\\{\Pi \times Y}\ar[r]&{Y}}
$$
correspondant est commutatif.

\begin{remarque}
\label{V.5.11}
On se gardera de confondre un pro-objet fondamental $P$ (qui n'est pas
muni d'une structure de groupe, et est connexe) avec le pro-groupe
fondamental (qui est un pro-\emph{groupe}, et en g\'en\'eral non
connexe). De fa\c con pr\'ecise, $G$ est connexe si
et
seulement si $\pi_F$ op\'erant sur lui-m\^eme par automorphismes
int\'erieurs est transitif, \ie si $\pi$ est r\'eduit \`a
l'\'el\'ement neutre, ou encore $\cal{C}$ \'equivalente \`a la
cat\'egorie des ensembles finis. Une autre diff\'erence
essentielle est que $G$ est d\'etermin\'e \`a isomorphisme
unique pr\`es, et $P$ n'est d\'etermin\'e qu'\`a isomorphisme
non unique pr\`es.

Soit $E$ un ensemble fini et consid\'erons le foncteur constant sur
le groupo\"ide~$\Gamma$, de valeur $E$: il d\'efinit en vertu
de~\ref{V.5.8} un objet de~$\cal{C}$, not\'e $E_{\cal{C}}$, et qui
peut aussi s'interpr\'eter comme la somme de~$E$ exemplaires de
l'objet final $e_{\cal{C}}$ de~$\cal{C}$. On peut consid\'erer
$E_{\cal{C}}$ comme un foncteur en $E$, de la cat\'egorie des
ensembles finis dans la cat\'egorie $\cal{C}$, et ce foncteur est
\emph{exact}, donc transforme groupes finis en $\cal{C}$-groupes,
etc. Si donc $X$ est un objet de~$\cal{C}$ sur lequel le groupe fini
$G$ op\`ere \`a droite, on voit qu'on peut consid\'erer $X$
comme un objet de~$\cal{C}$ ayant un $\cal{C}$-groupe d'op\'erateurs
\`a droite~$G_{\cal{C}}$. On dira donc par extension d'une
terminologie g\'en\'erale relative \`a des objets \`a
$\cal{C}$-groupes d'op\'erateurs, que $X$ est \emph{formellement
principal homogène}
\index{formellement principal homog\`ene (objet)|hyperpage}sous $G$ si $X$ est formellement principal homog\`ene sous
$G_{\cal{C}}$, \ie si le morphisme canonique
$$
X\times G_{\cal{C}} \to X\times X
$$
d\'eduit de l'op\'eration de~$G_{\cal{C}}$ sur~$X$ \`a droite,
est un isomorphisme. On dit que $X$ est \emph{principal homogène}
\index{principal homog\`ene (objet)|hyperpage}sous $G$ s'il l'est sous $G_{\cal{C}}$, \ie s'il l'est formellement,
et si de plus le quotient $X/G=X/G_{\cal{C}}$ est~$e_{\cal{C}}$. Si on
se fixe un foncteur fondamental, d'o\`u une \'equivalence de
$\cal{C}$ avec une cat\'egorie $\cal{C}(\pi)$, $X$ correspond \`a
un ensemble sur lequel $\pi$ op\`ere \`a gauche contin\^ument,
soit $E=F(X)$. Faire op\'erer $G$ sur~$X$ \`a droite revient alors
\`a faire op\'erer $G$ sur l'ensemble $E$ \`a droite, de fa\c con que les op\'erations de~$G$ commutent \`a celles de~$\pi$. On
constate alors aussit\^ot que $X$ est principal homog\`ene sous
$G$ si et seulement si l'ensemble $E$ est un espace principal
homog\`ene
sous $G$,
\ie si et seulement si $G$ y op\`ere de fa\c con simplement transitive. (D'ailleurs $X$ est formellement
principal homog\`ene \sss
$E$ principal homog\`ene \emph{ou} vide). Comparant
avec~\ref{V.5.3}, on voit que si $X$ est principal homog\`ene sous
$G$ \emph{et connexe}, alors l'homomorphisme donn\'e de~$G$ dans le
groupe oppos\'e \`a $\Aut(X)$ est un \emph{isomorphisme}; et
d'ailleurs pour qu'un objet $X$
de~$\cal{C}$
soit connexe et principal homog\`ene sous le groupe oppos\'e \`a
$\Aut(X)$, il faut et il suffit, avec les notations du \No \ref{V.4}, qu'il
soit isomorphe \`a un $P_i$ \emph{galoisien}. Dans le cas type
$\cal{C}=\cal{C}(\pi)$, cela signifie que $X$ est isomorphe \`a un
quotient de~$\pi$ par un sous-groupe invariant.

Supposons toujours donn\'e un foncteur fondamental~$F$. Alors la
donn\'ee d'un
$X$ principal homog\`ene sous un groupe fini $G$ op\'erant \`a
droite, et d'un point $a\in F(X)$, est \'equivalente \`a la
donn\'ee d'un homomorphisme de~$\pi$ dans le groupe~$G$. En effet,
\`a un tel homomorphisme on fait correspondre l'ensemble $E=G$, en y
faisant op\'erer~$\pi$ \`a gauche gr\^ace \`a l'homomorphisme
donn\'e $\pi\to G$ et les translations \`a gauche de~$G$, et en y
faisant op\'erer $G$ \`a droite par translation \`a droite, le
point marqu\'e $a$ de~$E$ \'etant l'\'el\'ement unit\'e
de~$G$. Gr\^ace \`a ce qui pr\'ec\`ede, on obtient ainsi de
fa\c con essentiellement unique tout triple $(X, G, a)$ ayant les
propri\'et\'es envisag\'ees plus haut,
puisqu'un
ensemble \`a
point marqu\'e principal homog\`ene sous un groupe $G$ s'identifie
\`a ce dernier. De cette fa\c con, on a une interpr\'etation
g\'eom\'etrique directe du foncteur $G\mto \Hom(\pi, G)$ de la
cat\'egorie des groupes finis dans~$\Ens$, foncteur qui est
pro-repr\'esentable \`a l'aide de~$\pi$, et dont la
consid\'eration fournirait donc une autre construction du groupe
$\pi$ associ\'e \`a~$F$.
\end{remarque}

\section{Foncteurs exacts d'une cat\'egorie galoisienne dans une autre}
\label{V.6}

\begin{proposition}
\label{V.6.1}
Soient $\cal{C}$, $\cal{C'}$ deux cat\'egories galoisiennes,
$H\colon \cal{C}\to \cal{C'}$ un foncteur covariant, $F'$ un foncteur
fondamental sur~$\cal{C'}$ et $F=F' \circ H$. Conditions
\'equivalentes:
\begin{enumerate}
\item[(i)] $H$ est \emph{exact}, \ie exact \`a gauche et exact
\`a droite.
\item[(ii)] $H$ est exact \`a gauche, transforme sommes finies en
sommes finies, et \'epimorphismes en \'epimorphismes (ou encore:
objets $\not\approx \emptyset_{\cal{C}}$ en objets $\not\approx
\emptyset_{\cal{C'}}$).
\item[(iii)] $F$ est un foncteur fondamental sur~$\cal{C}$.
\end{enumerate}
\end{proposition}

L'implication (i)$\To$(ii) est un fait g\'en\'eral aux
cat\'egories. D'ailleurs la premi\`ere forme donn\'ee \`a (ii)
implique la seconde, comme on voit en notant que si $X$ est un objet
de~$\cal{C}$, alors $X$ est $\not\approx \emptyset_{\cal{C}}$ \sss le
morphisme $X\to e_{\cal{C}}$ est un \'epimorphisme; on notera que
$F$ \'etant suppos\'e exact \`a gauche transforme $e_{\cal{C}}$
en~$e_{\cal{C'}}$. La deuxi\`eme forme donn\'ee \`a (ii)
implique~(iii), car $F$ \'etant exact \`a gauche donc
pro-repr\'esentable est justiciable de~\ref{V.5.6}
crit\`ere~(iii). Enfin (iii) implique~(i), comme il r\'esulte du
fait que $F$ est exact, et \og conservatif\fg (\ie satisfait l'axiome
(G~6) de~\no \ref{V.4}).

Soit alors $\Gamma$ le groupo\"ide fondamental de~$\cal{C}$,
$\Gamma'$ celui de~$\cal{C'}$. Donc si
$H$ est exact, alors $F' \mto F' \circ H$ est un foncteur du
groupo\"ide~$\Gamma'$ dans le groupo\"ide~$\Gamma$, qu'on
d\'enotera par $\tH$:
$$
\tH (F')(X)=F' (H(X))
$$
qu'on peut aussi \'ecrire, avec la notation $F(X)=E_X (F)$
introduite dans~\no \ref{V.6}:
$$
E_{H(X)} (F')=E_X (\tH(F')).
$$
Cette derni\`ere formule montre, compte tenu de~\ref{V.5.8}
ou~\ref{V.4.1}, que le foncteur exact $H$ est d\'etermin\'e (\`a
isomorphisme unique pr\`es) quand on conna\^it le foncteur
correspondant~$\tH$. Fixons nous un~$F'$, soit $F=\tH (F')$,
alors $\tH$ d\'efinit un homomorphisme de~$\pi_{F'} =\Aut(F')$
dans $\pi_F =\Aut(F)$:
$$
\tH\colon \pi_{F'} \to \pi_F \quad (F=\tH (F')=F' \circ H).
$$
D'ailleurs la formule plus haut montre (compte tenu de~\ref{V.5.8})
que cet homomorphisme a la propri\'et\'e suivante: pour tout
ensemble fini $E$ o\`u $\pi_F$ op\`ere contin\^ument, le groupe
$\pi_{F'}$ op\`ere \'egalement \emph{continûment} gr\^ace
\`a l'homomorphisme pr\'ec\'edent $\pi_{F'} \to
\pi_F$. \hbox{Appliquant} ceci aux quotients de~$\pi_F$ par ses sous-groupes
invariants ouverts, on voit que la condition pr\'ec\'edente
signifie aussi que l'homomorphisme consid\'er\'e est
continu. R\'eciproquement,
donnons-nous
un objet $F$ de~$\Gamma$, un
objet $F'$ de~$\Gamma'$ et un homomorphisme continu
$$
u\colon \pi_{F'} \to \pi_F ,
$$
il lui correspond donc un foncteur de~$\cal{C} (\pi)$
dans~$\cal{C}(\pi')$, manifestement exact, donc en vertu
de~\ref{V.4.1} il lui correspond un foncteur $H$ de~$\cal{C}$ dans
$\cal{C'}$ qui est exact, et tel que $\tH\colon \pi_{F'} \to \pi_F$
soit pr\'ecis\'ement~$u$. On peut aussi, au lieu d'un
homomorphisme de groupes, partir d'un \emph{foncteur}
$$
U\colon \Gamma' \to \Gamma
$$
qui est tel que pour \emph{tout} $F' \in \Gamma'$ (ou \emph{un} $F'
\in \Gamma'$, cela revient au m\^eme) l'homomorphisme correspondant
$\pi_{F'} \to \pi_F$ soit continu: un tel foncteur est isomorphe \`a
un foncteur de la forme~$\tH$, o\`u $H\colon \cal{C} \to
\cal{C'}$ est un foncteur exact d\'etermin\'e \`a isomorphisme
unique pr\`es. Ainsi:

\begin{corollaire}
\label{V.6.2}
Pour
qu'un foncteur $H\colon \cal{C} \to \cal{C'}$ de cat\'egories
galoisiennes soit exact, il faut et il suffit qu'il existe des
\'equivalences $\cal{C}(\pi)\to\cal{C}$ et
$\cal{C'}\to\cal{C}(\pi')$
qui transforment le foncteur $H$ en le foncteur
$\cal{C}(\pi)\to\cal{C}(\pi')$ associ\'e \`a un homomorphisme de
groupes topologiques $\pi'\to\pi$.
\end{corollaire}

\begin{corollaire}
\label{V.6.3}
Soient $\cal{C}$, $\cal{C'}$ deux cat\'egories galoisiennes,
$\Gamma$, $\Gamma'$ leurs groupo\"ides fondamentaux. Alors la
cat\'egorie des foncteurs exacts de~$\cal{C}$ dans $\cal{C'}$ est
\'equivalente \`a la cat\'egorie des foncteurs $U\colon \Gamma'
\to \Gamma$ ayant la propri\'et\'e suivante: pour tout $F'$ dans~$\Gamma'$ (ou \emph{un} $F'$ dans~$\Gamma'$, cela revient au
m\^eme), posant $F=U(F')$, l'homomorphisme
$$
\pi_{F'} =\Aut (F')\to \pi_F =\Aut (F)
$$
d\'efini par $U$ est continu.
\end{corollaire}

Consid\'erons le pro-groupe fondamental $\Pi$ de~$\cal{C}$, alors un
foncteur exact $H$ le transforme en un pro-groupe $H(\Pi)$
de~$\cal{C'}$, nous allons d\'efinir un homomorphisme
$$
\Pi' \to H(\Pi)
$$
(o\`u $\Pi'$ est le pro-groupe fondamental de~$\cal{C'}$), par la
condition que pour tout objet $F'$ de~$\Gamma'$, l'homomorphisme
correspondant
$$
F' (\Pi')=\pi_{F'} \to F' (H(\Pi))=\pi_F \qquad
\text{(où $F=F'\circ H=\tH(F')$)}
$$
soit l'homomorphisme naturel
$$
\Aut(F')\to \Aut(F' \circ H).
$$
Comme ce dernier est fonctoriel en~$F'$, il d\'efinit bien en vertu
de~\ref{V.5.8} un homomorphisme de pro-objets, et en fait de
pro-groupes, de~$\cal{C'}$. Cet homomorphisme est dit
\emph{associé} au foncteur~$H$.

Soit maintenant $H'$ un deuxi\`eme foncteur exact, de la
cat\'egorie galoisienne $\cal{C'}$ dans une cat\'egorie
galoisienne~$\cal{C''}$. Il est trivial qu'on a
$$
{}^t (H' H)=\tH\,\tH'
$$
(N.B. on a l\`a une identit\'e de foncteurs, et non seulement un
isomorphisme canonique). On a une propri\'et\'e de
transitivit\'e analogue pour les homomorphismes associ\'es des
pro-groupes fondamentaux.

Nous
allons maintenant interpr\'eter les propri\'et\'es du foncteur
exact $H$ en termes de l'homomorphisme correspondant
$$
u\colon \pi_{F'} \to \pi_{F} \qquad \text{(où $F' =F'\circ H$)}.
$$
Il est commode d'introduire la notion d'\emph{objet ponctué}
\index{ponctu\'e (objet)|hyperpage}de la cat\'egorie galoisienne $\cal{C}$ (muni de son foncteur
fondamental $F$): c'est par d\'efinition un objet $X$ de~$\cal{C}$
muni d'un \'el\'ement $a$ de~$F(X)$. Il s'interpr\`ete donc
comme un ensemble fini o\`u $\pi_F$ op\`ere contin\^ument \`a
gauche, muni d'un point~$a$. Donc les objets ponctu\'es
\emph{connexes} de~$\cal{C}$ s'identifient en vertu de~\ref{V.5.3} aux
sous-groupes ouverts de~$\pi_F$. Si $U$ et $V$ sont deux tels
sous-groupes, correspondants \`a des objets ponctu\'es connexes
$X$, $Y$ de~$\cal{C}$, alors il existe un homomorphisme ponctu\'e de
$X$ dans $Y$ si et seulement si $U\subset V$, et cet homomorphisme est
alors unique. Bien entendu, le foncteur $H$ transforme objets
ponctu\'es en objets ponctu\'es (puisque $F=F' \circ H$). Notons
d'autre part qu'un sous-groupe ferm\'e d'un groupe tel que $\pi_F$
est l'intersection des sous-groupes ouverts qui le contiennent; par
suite,
si
$M$, $N$ sont deux sous-groupes ferm\'es, alors $M\subset N$ si et
seulement si tout sous-groupe ouvert qui contient $N$ contient
\'egalement~$M$. Gr\^ace \`a ces remarques, on prouve facilement
les r\'esultats qui suivent:

\begin{proposition}
\label{V.6.4}
Soit $X$ un objet ponctu\'e connexe de~$\cal{C}$, associ\'e \`a
un sous-groupe ouvert $U$ de~$\pi_F$. Pour que
$U$
contienne $u(\pi_{F'})$ (\resp le sous-groupe invariant ferm\'e
engendr\'e par~$u(\pi_{F'})$) il faut et il suffit que $H(X)$
admette une section ponctu\'ee (\resp soit compl\`etement
d\'ecompos\'e).
\end{proposition}

On appellera \emph{section} --- sous-entendu: au-dessus de l'objet
final --- d'un objet $X$ d'une cat\'egorie galoisienne~$\cal{C}$, un
morphisme de l'objet final $e_{\cal{C}}$ dans~$X$, ce qui revient
\`a la donn\'ee d'un \'el\'ement $a$ de~$F(X)$ invariant par
$\pi_F$; si $X$ est ponctu\'e, on dit qu'on a une \emph{section
ponctuée} si elle est compatible avec les structures ponctu\'ees
sur~$X$ et~$e_{\cal{C}}$, \ie si $a$ est pr\'ecis\'ement l'objet
marqu\'e de~$F(X)$. Une telle section est donc unique, et existe si
et seulement si l'objet marqu\'e de~$F(X)$ est invariant
par~$\pi_F$. Enfin, un objet d'une cat\'egorie galoisienne est dit
\emph{complètement décomposé}
\index{compl\`etement d\'ecompos\'e (objet)|hyperpage}\index{decompose (objet completement)@d\'ecompos\'e (objet compl\`etement)|hyperpage}s'il est isomorphe \`a une somme d'objets finaux, \ie si $\pi_F$
op\`ere trivialement dans $F(X)$ --- condition \'evidemment plus
forte que l'existence d'une section ponctu\'ee, lorsque $X$ est
ponctu\'e. La proposition~\ref{V.6.4} r\'esulte trivialement des
d\'efinitions et remarques pr\'ec\'edentes.

\begin{corollaire}
\label{V.6.5}
Pour
que $u$ soit trivial, il faut et il suffit que pour tout objet
$X$ de~$\cal{C}$, $H(X)$ soit compl\`etement d\'ecompos\'e.
\end{corollaire}

\begin{proposition}
\label{V.6.6}
Soit $X'$ un objet ponctu\'e connexe de~${\cal{C}}'$, associ\'e
\`a un sous-groupe ouvert $U'$ de~$\pi_{F'}$. Pour que $U'$
contienne~$\Ker u$, il faut et il suffit qu'il existe un objet
ponctu\'e connexe $X$ de~$\cal{C}$ et un homomorphisme ponctu\'e
de la composante connexe ponctu\'ee $X'_0$ de~$H(X)$ dans~$X'$ (donc
que $X$ soit isomorphe comme objet ponctu\'e \`a un quotient de la
composante connexe neutre de l'image inverse d'un objet ponctu\'e
de~$\cal{C}$). Si $u$ est surjectif, la condition pr\'ec\'edente
\'equivaut aussi \`a la suivante: $X$ est isomorphe \`a
un~$H(X)$, o\`u $X$ est un objet ponctu\'e de~$\cal{C}$.
\end{proposition}

(On appelle \emph{composante connexe neutre}
\index{neutre (composante connexe d'un objet ponctu\'e)|hyperpage}d'un objet ponctu\'e $X$ d'une cat\'egorie galoisienne~$\cal{C}$,
l'unique sous-objet connexe ponctu\'e de~$X$; il correspond \`a la
trajectoire sous $\pi_F$ du point marqu\'e de~$F(X)$, en vertu
de~\ref{V.5.3}). Comme le fait que $U'$ contienne $\Ker u$ ne
d\'epend pas de la ponctuation choisie de~$X'$ (car une autre
ponctuation revient \`a remplacer $U$ par un sous-groupe
conjugu\'e \`a~$U$), on voit:

\begin{corollaire}
\label{V.6.7}
Pour que $U'$ contienne~$\Ker u$, il faut et il suffit qu'il existe un
\hbox{objet~$X$} de~$\cal{C}$ (qu'on peut supposer connexe) et un morphisme
d'une composante connexe de~$H(X)$ dans~$X'$. Si $u$ est surjectif,
cela signifie aussi que $X'$ est isomorphe \`a un objet de la
forme~$H(X)$.
\end{corollaire}

\begin{corollaire}
\label{V.6.8}
Pour que $u$ soit injectif, il faut et il suffit que pour tout objet
$X'$ de~${\cal{C}}'$, il existe un objet $X$ de~$\cal{C}$ et un
homomorphisme d'une composante connexe de~$H(X)$ dans~$X'$.
\end{corollaire}

\begin{proposition}
\label{V.6.9}
Les conditions suivantes sont \'equivalentes:
\begin{enumerate}
\item[(i)] L'homomorphisme $u\colon \pi_{F'} \to \pi_F$ est surjectif.
\item[(ii)] Pour tout objet connexe $X$ de~$\cal{C}$, $H(X)$ est
connexe.
\item[(iii)] Le foncteur $H$ est pleinement fid\`ele.
\end{enumerate}
\end{proposition}

Ce dernier fait signifie que pour deux objets~$X$, $Y$ de~$\cal{C}$,
l'application naturelle
$$
\Hom(X,Y) \to \Hom(H(X),H(Y))
$$
est bijective.

\begin{corollaire}
\label{V.6.10}
Pour
que $u$ soit un isomorphisme, il faut et suffit que $H$ soit une
\'equivalence de cat\'egories, ou encore que les deux conditions
suivantes soient v\'erifi\'ees:
\begin{enumerate}
\item[a)] pour tout objet connexe $X$ de~$\cal{C}$, $H(X)$ est connexe
\item[b)] tout objet de~$\cal{C}'$ est isomorphe \`a un objet de la
forme~$H(X)$.
\end{enumerate}
\end{corollaire}

\begin{proposition}
\label{V.6.11}
Soient $H\colon \cal{C}\to \cal{C}'$ et $H'\colon\cal{C}\to \cal{C}''$
des foncteurs exacts entre cat\'egories galoisiennes, $F''$ un
foncteur fondamental sur~$\cal{C}''$, posons $F'=F''H'$ et $F=F'H$, et
consid\'erons les homomorphismes associ\'es
$$
u'\colon\pi_{F''}\to \pi_{F'}\qquad u\colon\pi_{F'}\to \pi_F.
$$
Pour que $\Ker u\subset \Im u'$ \ie pour que $uu'$ soit
l'homomorphisme trivial, il faut et il suffit que pour tout objet $X$
de~$\cal{C}$, $H'(H(X))$ soit compl\`etement d\'ecompos\'e. Pour
que $\Ker u \supset \Im u' $, il faut et il suffit que pour tout objet
ponctu\'e connexe $X'$ de~$\cal{C}'$ tel que $H'(X')$ admette une
section ponctu\'ee, il existe un objet~$X$ de~$\cal{C}$ et un
homomorphisme d'une composante connexe de~$H(X)$ dans~$X'$.
\end{proposition}

La premi\`ere assertion r\'esulte de la derni\`ere affirmation
de~\ref{V.6.4}. La deuxi\`eme r\'esulte de la conjonction
de~\ref{V.6.4} et~\ref{V.6.6}.

\begin{remarque}
\label{V.6.12}
Il n'est pas vrai en g\'en\'eral, sous les conditions
de~\ref{V.6.8} que $X'$ soit isomorphe \`a un objet de la forme
$H(X)$. On peut montrer que pour que tout objet connexe (donc tout
objet) de~$\cal{C}'$ soit isomorphe \`a un objet de la forme $H(X)$,
il faut et il suffit que $u$ soit un isomorphisme de~$\pi_{F'}$ sur un
sous-groupe \emph{facteur direct} de~$\pi_F$. En pratique cependant,
on construit directement un homomorphisme $\pi_F\to\pi_{F'}$ inverse
\`a droite de~$u$ \`a l'aide d'un foncteur exact convenable
de~$\cal{C'}$ dans~$\cal{C}$.
\end{remarque}

\begin{proposition}
\label{V.6.13}
Soient $\cal{C}$ une cat\'egorie galoisienne munie d'un foncteur
fondamental $F$, $S$ un objet connexe de~$\cal{C}$, $\cal{C}'$ la
cat\'egorie des objets de~$\cal{C}$ au-dessus de~$S$. Alors
$\cal{C'}$ est une cat\'egorie galoisienne, et le foncteur $X\mto
H(X)=X\times S$ de~$\cal{C}$ dans~$\cal{C}'$ est exact. Soit $a\in
F(S)$, et soit $F'$ le foncteur de~$\cal{C}'$ dans la cat\'egorie
des ensembles finis d\'efini par
$$
F'(X')=\text{image inverse de~$a$ par }F(X')\to F(S).
$$
Alors on a un isomorphisme $F\cong F'\circ H$, et l'homomorphisme
correspondant
$$
u:\pi_{F'}\to \pi_F
$$
est un isomorphisme de~$\pi_{F'}$ sur le sous-groupe ouvert $U$ de
$\pi_{F}$ stabilisateur de l'\'el\'ement marqu\'e $a$ de~$F(X)$.
\end{proposition}

La d\'emonstration est laiss\'ee au lecteur.

\section{Cas des pr\'esch\'emas}
\label{V.7}
Soit $S$ un pr\'esch\'ema localement noeth\'erien et
\emph{connexe}, et soit
$$
a\colon \Spec(\Omega)\to S
$$
un point g\'eom\'etrique de~$S$, \`a valeurs dans un corps
alg\'ebriquement clos $\Omega$. On posera
$$
\cal{C}=\text{catégorie des revêtements étales de~$S$}
$$
et pour un objet $X$ de~$\cal{C}$, \ie un rev\^etement \'etale
$X$ de~$S$, on pose
$$
F(X)=\text{ensemble des points géométriques de~$X$ au-dessus
de~$a$.}
$$
Ainsi, $F$ devient un foncteur sur~$\cal{C}$ \`a valeurs dans la
cat\'egorie des ensembles finis. Les propri\'et\'es (G~1) \`a
(G~6) sont satisfaites: pour (G~1), c'est contenu dans les sorites de
I~\ref{I.4.6}, (G~2) r\'esulte de~\ref{V.3.4}, (G~3) de~\ref{V.3.5},
(G~4) est trivial par d\'efinition, (G~5) r\'esulte de~\ref{V.3.5} et
du d\'ebut du \No \ref{V.2}, enfin (G~6) est prouv\'e dans~\ref{V.3.7}. On
peut donc appliquer les r\'esultats des \No \ref{V.4}, \ref{V.5}, \ref{V.6}. Cela permet en
particulier de d\'efinir un pro-objet $P$ de~$\cal{C}$
repr\'esentant $F$, appel\'e \emph{revêtement universel de~$S$
au point~$a$},
\index{revetement universel d'un preschema en un point@rev\^etement universel d'un pr\'esch\'ema en un point|hyperpage}et un groupe topologique $\pi=\Aut(F)=\Aut^0 P$, appel\'e
\emph{groupe fondamental de~$S$ en~$a$},
\index{groupe fondamental d'un pr\'esch\'ema en un point|hyperpage}et not\'e~$\pi_1(S,a)$.
\label{indnot:eg}Le foncteur $F$ d\'efinit alors une \'equivalence de la
cat\'egorie $\cal{C}$ avec la cat\'egorie des ensembles finis
o\`u $\pi=\pi_1(S,a)$ op\`ere contin\^ument. Cette
\'equivalence permet donc d'interpr\'eter les op\'erations
courantes de limites projectives et limites inductives finies sur des
rev\^etements (produits, produits fibr\'es, sommes, passage au
quotient, etc.) en termes des op\'erations analogues dans
$\cal{C}(\pi)$, \ie en termes des op\'erations \'evidentes sur
des ensembles finis o\`u $\pi$ op\`ere. Notons d'ailleurs, puisque
les composantes connexes topologiques d'un rev\^etement \'etale
sont \'egalement des rev\^etements \'etales, qu'\emph{un \hbox{objet $X$} de~$\cal{C}$ est connexe dans $\cal{C}$ si et seulement si il est
topologiquement connexe}; en vertu de~\ref{V.5.3} cela signifie donc
que $\pi_1$ op\`ere transitivement dans~$F(X)$.
Notons que pour qu'un objet $X$ de~$\cal{C}$ soit fid\`element plat
et quasi-compact sur~$S$ (comme il est d\'ej\`a plat et
quasi-compact sur~$S$), il faut et suffit que $X\to S$ soit surjectif
\ie soit un \'epimorphisme dans $\cal{C}$, ou encore que $X\neq
\emptyset$. On conclut alors du crit\`ere \ref{V.2.6}~(iii) que~$X$
\emph{est un revêtement principal de
$S$
de groupe $G$ si et seulement si il est un espace principal
homogène sous $G$ dans la catégorie $\cal{C}$}, (tel qu'il a
\'et\'e d\'efini dans \No \ref{V.5}).

Si $a'$ est un autre point g\'eom\'etrique de~$S$ (correspondant
\`a un corps alg\'ebriquement clos $\Omega'$, qui peut \^etre
diff\'erent de~$\Omega$ et qui peut m\^eme avoir une
caract\'eristique diff\'erente), il d\'efinit un foncteur fibre
$F'=F_{a'}$ de~$\cal{C}$ dans la cat\'egorie des ensembles finis,
qui est encore exact, donc isomorphe \`a $F=F_a$. Par suite, les
groupes fondamentaux $\pi_1(S;a)$
\label{indnot:eg1}pour $a$ variable sont isomorphes entre eux. Si on d\'esigne par
$\pi_1(S;a,a')$
\label{indnot:eh}l'ensemble des isomorphismes (ou ce qui revient au m\^eme,
l'ensembles des homomorphismes) $F_a\to F_{a'}$ des foncteurs fibres
associ\'es, on obtient ainsi un \emph{groupoïde} dont
l'ensemble des objets est l'ensemble des points g\'eom\'etriques
de~$S$, les groupes fondamentaux \'etant les groupes
d'automorphismes des objets dudit groupo\"ide. L'ensemble
$\pi_1(S;a',a)$ peut \^etre appel\'e l'\emph{ensemble des classes
de chemins de~$a$ à $a'$}. Ces classes se composent donc de fa\c con \'evidente. Enfin, on peut d\'efinir un pro-groupe $\Pi_1^S$
de~$\cal{C}$, qu'on pourra appeler \emph{pro-groupe fondamental
de}~$S$
\index{pro-groupe fondamental d'un pr\'esch\'ema|hyperpage}ou \emph{système local des groupes fondamentaux sur}~$S$,
\index{systeme local des groupes fondamentaux@syst\`eme local des groupes fondamentaux|hyperpage}d\'efini \`a isomorphisme unique pr\`es par la condition qu'on
ait un isomorphisme, fonctoriel en le point g\'eom\'etrique $a$ de
$S$:
$$
F_a(\Pi_1^S)=\pi_1(S;a)
$$
(\cf remarque~\ref{V.5.10}). En particulier, si $s$ est un point
ordinaire de~$S$, la fibre de~$G$ en~$s$ est un pro-groupe sur~$\kres(s)$,
limite projective de groupes finis \'etales sur~$\kres(s)$; on pourrait
appeler ce pro-groupe \emph{le groupe fondamental de~$S$ en le point
ordinaire $s$ de~$S$}, et le noter $\pi_1(S,s)$. Par d\'efinition, ces
points \`a valeurs dans une extension alg\'ebriquement close
$\Omega$ de~$\kres(s)$ sont les \'el\'ements de~$\pi_1(S;a)$, o\`u
$a$ est le point g\'eom\'etrique de~$S$ d\'efini par ladite
extension. En particulier, (prenant pour $S$ le spectre d'un corps)
\`a tout corps $k$ est associ\'e canoniquement et fonctoriellement
un pro-groupe sur~$k$, qu'on pourrait noter $\pi_1(k)$, limite
projective de groupes finis \'etales sur~$k$, et dont les points
dans une extension alg\'ebriquement close $\Omega$ de~$k$
s'identifient
aux \'el\'ements du groupe de Galois topologique de~$\bar{k}/k$,
o\`u $\bar{k}$ est la cl\^oture galoisienne de~$k$ dans $\Omega$
(\cf \ref{V.8.1}). Ce groupe $\pi_1(k)$ ne semble pas encore avoir
retenu l'attention des alg\'ebristes.

Soit maintenant
$$
f\colon S'\to S
$$
un morphisme d'un pr\'esch\'ema connexe localement noeth\'erien
dans un autre, soit $a'$ un point g\'eom\'etrique de~$S'$ et soit
$a=f(a')$ son image directe dans $S$. Alors le foncteur \og image
inverse\fg induit un foncteur de la cat\'egorie $\cal{C}(S)$
\label{indnot:ei}des rev\^etements \'etales de~$S$, dans la cat\'egorie
$\cal{C}(S')$ des rev\^etements \'etales de~$S'$:
$$
f^\bbullet:\cal{C}(S)\to \cal{C}(S')
$$
On a d'ailleurs un isomorphisme de foncteurs
$$
F_a\cong F_{a'}\circ f^\bbullet,
$$
de sorte que $f^\bbullet$ est un foncteur \emph{exact}, auquel
s'appliquent les r\'esultats du \No \ref{V.6}. On a en particulier un
homomorphisme canonique
$$
u=\pi_1(f;a'):\pi_1(S',a')\to \pi_1(S;a) \qquad (a'=f(a))
$$
\label{indnot:ej}qui permet de reconstituer le foncteur image inverse, comme une
op\'eration de restriction des groupes d'op\'erateurs. Les
propri\'et\'es du foncteur $f^\bbullet$ s'expriment donc de
fa\c con simple par les propri\'et\'es de l'homomorphisme de
groupes associ\'es, comme il a \'et\'e explicit\'e dans le
\No \ref{V.6}. Si en particulier $S'$ est un rev\^etement \'etale de~$S$,
alors l'homomorphisme $u$ est un isomorphisme de~$\pi_1(S',a')$ sur le
sous-groupe ouvert de~$\pi_1(S,a)$ qui d\'efinit le rev\^etement
\'etale connexe ponctu\'e $S'$ de~$S$ (\ie le stabilisateur~$U$
de~$a'\in F_a(S')$ dans $\pi_1(S;A)$).

Si on d\'esire interpr\'eter les homomorphismes $\pi_1(f;a')$ pour
un point g\'eom\'etrique variable $a'$, on doit, conform\'ement
\`a ce qui a \'et\'e dit dans le \No \ref{V.6}, consid\'erer un
homomorphisme
$$
\Pi_1(f):\Pi_1^{S'}\to f^\bbullet(\Pi_1^S)
$$
de pro-groupes sur~$S$, et prendre l'homomorphisme correspondant pour
les fibres g\'eom\'etriques.

\section{Cas d'un pr\'esch\'ema de base normale}
\label{V.8}
\begin{proposition}
\label{V.8.1}
Soit $S$ le spectre d'un corps $k$, et soit $\Omega$ une extension
alg\'ebriquement close de~$k$, d\'efinissant un point
g\'eom\'etrique $a$ de~$S$ \`a valeurs dans $\Omega$. Soit
$\bar{k}$ la cl\^oture s\'eparable de~$k$ dans $\Omega$. Alors il
existe un isomorphisme canonique de~$\pi_1(S,a)$ sur le groupe de
Galois topologique de~$\bar{k}/k$.
\end{proposition}

Soit $k'$ la cl\^oture alg\'ebrique de~$k$ dans $\Omega$, il
correspond donc \`a un point g\'eom\'etrique~$b$ de~$S$, \`a
valeurs dans $k'$. L'homomorphisme naturel de foncteurs $F_b\to F_a$
est \'evidemment un isomorphisme, car un $k$-homomorphisme d'une
extension finie s\'eparable de~$k$ dans $\Omega$ prend
n\'ecessairement ses valeurs dans $\bar{k}$ et a fortiori dans
$k'$. D'autre part, le groupe $\pi'$ des $k$-automorphismes de~$k'/k$
op\`ere de fa\c con \'evidente sur~$F_b$, d'o\`u un
homomorphisme $\pi'\to \Aut(F_b)\isomto \Aut(F_a)=\pi_1(S;a)$.
D'autre part, il est bien connu que l'homomorphisme naturel de~$\pi'$
dans le groupe $\pi$ des automorphismes de~$\bar{k}/k$ est un
isomorphisme. On obtient ainsi un homomorphisme canonique $\pi\to
\pi_1(S;a)$, reste \`a montrer que c'est un isomorphisme. En effet,
cet homomorphisme est injectif, car un \'el\'ement du noyau est un
automorphisme de~$\bar{k}/k$ qui induit l'identit\'e sur toute
sous-extension s\'eparable finie, donc est trivial. Cet
homomorphisme est surjectif, car si $X$ est un rev\^etement
\'etale \emph{connexe} de~$S$, donc d\'efini par une
\emph{extension} finie s\'eparable~$L/k$, alors $\pi$ est transitif
sur l'ensemble des $k$-homomorphismes de~$L$ dans $k'$, comme bien
connu.

\begin{proposition}
\label{V.8.2}
Soient $S$ un pr\'esch\'ema connexe, localement noeth\'erien et
normal, $K=\kres(s)$ son corps des fonctions $=$ le corps r\'esiduel en
son point g\'en\'erique $s$, $\Omega$ une extension
alg\'ebriquement close de~$K$, d\'efinissant un point
g\'eom\'etrique $a'$ de~$S'=\Spec(K)$ et un point
g\'eom\'etrique $a$ de~$S$. Alors l'homomorphisme
$\pi_1(S';a')\to\pi_1(S;a)$ est surjectif. Lorsqu'on identifie le
premier groupe au groupe de Galois de la cl\^oture s\'eparable
$\bar{K}$ de~$K$ dans~$\Omega$ (\cf \ref{V.8.1}) alors le noyau de
l'homomorphisme pr\'ec\'edent correspond par la th\'eorie de
Galois \`a la sous-extension de~$\bar{K}/K$ compos\'ee des
extensions finies de~$K$ dans $\Omega$ qui sont non ramifi\'ees sur~$S$.
\end{proposition}

La premi\`ere assertion signifie que l'image inverse sur~$S'$ d'un
rev\^etement \'etale connexe $X$ de~$S$ est connexe, \ie que $X$
est int\`egre, ce n'est autre que (I~\ref{I.10.1}). Le noyau de
l'homomorphisme pr\'ec\'edent s'interpr\`ete alors comme
form\'e des automorphismes de~$\bar{K}/K$ qui induisent
l'identit\'e sur les ensembles~$F_a(X)$,
o\`u on peut supposer le rev\^etement \'etale $X$ de~$S$
connexe. Mais cela signifie que cet automorphisme induit
l'identit\'e sur les sous-extensions finies de~$\bar{K}/K$ qui sont
non ramifi\'ees sur~$S$, ce qui prouve la derni\`ere assertion.

\begin{remarquestar}
Gr\^ace \`a cette interpr\'etation du groupe fondamental du
pr\'esch\'ema normal $S$ en termes de th\'eorie des Galois
habituelle, la d\'efinition \'etait connue dans ce cas depuis
longtemps.
\end{remarquestar}

\section{Cas des pr\'esch\'emas non connexes: cat\'egories
multigaloisiennes}
\label{V.9}

Soit $S$ un pr\'esch\'ema localement noeth\'erien, et soient
$(S_i)_{i\in I}$ ses composantes connexes. Alors la cat\'egorie
$\cal{C}(S)$ des rev\^etements \'etales de~$S$ est \'equivalente
\`a la cat\'egorie produit des $\cal{C}(S_i)$, qui
s'interpr\`etent en termes des groupes fondamentaux des $S_i$, une
fois choisie un point g\'eom\'etrique dans chaque $S_i$. Dans
l'application de la th\'eorie de la descente pour les morphismes
\'etales, il est parfois malcommode de faire choix pour tout $S_i$
d'un point g\'eom\'etrique de~$S_i$. Il est plus commode alors
de recourir \`a la g\'en\'eralisation naturelle de~\ref{V.5.8}
pour interpr\'eter $\cal{C}(S)$ comme une cat\'egorie de foncteurs
sur le groupo\"ide des points g\'eom\'etriques de~$S$,
consid\'er\'e comme somme des groupo\"ides correspondants aux
composantes connexes de~$S$; les foncteurs en question sont les
foncteurs \`a valeurs dans la cat\'egorie des ensembles finis,
satisfaisant la propri\'et\'e de continuit\'e analogue \`a
celle invoqu\'ee dans~\ref{V.5.8}. En pratique, on aura une famille
$(a_t)_{t\in E}$ de points g\'eom\'etriques de~$S$, telle que
toute composante connexe $S_i$ de~$S$ en contienne au moins un, et on
pourra alors, comme dans~\ref{V.5.10}, remplacer le groupo\"ide de
tous les points g\'eom\'etriques de~$S$ par le groupo\"ide
analogue dont l'ensemble sous-jacent est~$E$. Bien entendu, ces
consid\'erations devraient s'ins\'erer dans des d\'efinitions
g\'en\'erales concernant les cat\'egories qui sont
\'equivalentes \`a des cat\'egories produits de cat\'egories
de la forme $\cal{C}(\pi)$, et qu'on pourra appeler
\emph{catégories multi-galoisiennes}.
\index{multi-galoisienne (cat\'egorie)|hyperpage}Nous en laisserons le d\'etail au lecteur.

